% 第 6 章 标准形理论 (Normalformentheorie)
\chapter{标准形理论}\index{biaozhunxinglilun@标准形理论(Normalformtheorie)}

\section*{预备说明}

本章的目标是弄清有限维 $K$-向量空间 $V$ 的自同态的结构.但究竟应把自同态 $f \colon V \longrightarrow V$ 的"结构"理解为什么呢? 例如可以提出下面的问题:
\begin{enumerate}
\item 是否存在非平凡的线性子空间 $U \subset V$ 满足 $f(U) \subset U$, 使得 $f|_U$ 具有特别简单的形状, 例如 $f|_U = \lambda \operatorname{id}_U$, 其中 $\lambda \in K$ 是某个纯量?
\item 为了在整个 $V$ 上描述 $f$: 能否把 $V$ 分解为非平凡线性子空间的直和 $V = \bigoplus_{i = 1}^{r} U_i$, 使得 $f(U_i) \subset U_i$, 并且 $f|_{U_i}$ 能以有意义的方式刻画? 是否存在这种类型的最细分解, 它是否以某种方式被唯一地刻画?
\end{enumerate}

这些已经是我们想研究的最重要的问题.首先, 第一个问题的回答相对简单.对 $\lambda \in K$, 考虑 $V$ 的自同态 $f - \lambda \operatorname{id}$.它的核恰好指出 $V$ 中使 $f$ 表现得如同 $\lambda \operatorname{id}$ 的那个(极大的)线性子空间, 而这个子空间非平凡, 当且仅当 $f - \lambda \operatorname{id}$ 的核非平凡, 即按推论~\mtref{cor:2.1.11} 恰在 $f - \lambda \operatorname{id}$ 不可逆时, 从而按定理~\mtref{thm:4.3.3} 恰在 $\det(f - \lambda \operatorname{id}) = 0$ 时.因此要对 $\lambda \in K$ 解方程 $\det(f - \lambda \operatorname{id}) = 0$, 我们由此将自然而然地考虑 $f$ 的所谓\emph{特征多项式}(\textup{charakteristisches Polynom})\index{tezhengduoxiangshi@特征多项式(charakteristisches Polynom)} $\chi_f \in K[T]$, 它是在对 $\det(\lambda \operatorname{id} - f)$ 套用描述 $f$ 的矩阵的行列式定义时,把 $\lambda$ 换成变量 $T$ 而得到的.

$\chi_f$ 在 $K$ 中的零点称为 $f$ 的\emph{特征值}(\textup{Eigenwert})\index{tezhengzhi@特征值(Eigenwert)}.对这样的特征值 $\lambda$, 称 $V_\lambda = \ker(\lambda \operatorname{id} - f)$ 为属于 $\lambda$ 的\emph{特征空间}(\textup{Eigenraum})\index{tezhengkongjian@特征空间(Eigenraum)}, 而 $V_\lambda - \{0\}$ 中的元素称为属于特征值 $\lambda$ 的\emph{特征向量}(\textup{Eigenvektor})\index{tezhengxiangliang@特征向量(Eigenvektor)}.我们将证明, 属于不同特征值的特征向量总是线性无关的, 并由此推出 $f$ 的属于不同特征值的诸特征空间之和恒为直和.若这个和已经与 $V$ 一致, 则称 $f$ \emph{可对角化}(\textup{diagonalisierbar})\index{keduijiaohua@可对角化(diagonalisierbar)}, 我们以此表示 $V$ 有一个由 $f$ 的特征向量组成的基, 其结果是 $f$ 的描述矩阵为对角矩阵.例如当 $f$ 的不同特征值的个数等于 $V$ 的维数时, 就出现这种情况.我们也由此得到对开头问题 2 的第一个(部分的)回答.

比较容易看出, 线性映射 $f \colon V \longrightarrow V$ 一般不可对角化.例如 $\mathbb{R}^2$ 中绕 $0$ 的旋转不可对角化, 除非\emph{旋转角}(\textup{Drehwinkel})\index{xuanzhuanjiao@旋转角(Drehwinkel)}为 $0^{\circ}$ 或 $180^{\circ}$.于是为回答问题 2, 还必须寻求构造线性子空间 $U \subset V$(满足 $f(U) \subset U$, 即所谓 $f$-不变的线性子空间)的其他可能性.下面的观察是基础性的: 对向量 $x \in V$, 考虑由元素 $x, f(x), f^2(x), \ldots$ 生成的线性子空间 $U \subset V$.则 $U$ 显然是 $V$ 中一个 $f$-不变的线性子空间, 显然是包含 $x$ 的最小者.我们称 $U$ 为 $V$ 中由 $x$ 生成的 $f$-循环线性子空间.若利用第 5 章关于多项式环 $K[T]$ 的结果, 它的结构就容易描述.事实上可以考虑 $K$-线性映射

\[
\varphi \colon K[T] \longrightarrow U, \quad \sum_{i \in \mathbb{N}} c_i T^i \longmapsto \sum_{i \in \mathbb{N}} c_i f^i(x),\]

并注意到 $\ker \varphi$ 是 $K[T]$ 中的一个理想.因为对

\[
p = \sum_{i = 0}^{r} c_i T^i \in K[T], \qquad q = \sum_{j = 0}^{s} d_j T^j \in \ker \varphi\]

有

\[
\begin{array}{c}
\varphi(pq) = \varphi\bigg(\sum_{k = 0}^{r + s}\Big(\sum_{i + j = k} c_i d_j\Big) T^k\bigg)
= \sum_{k = 0}^{r + s}\Big(\sum_{i + j = k} c_i d_j\Big) f^k(x) \\
= \Big(\sum_{i = 0}^{r} c_i f^i\Big)\Big(\sum_{j = 0}^{s} d_j f^j\Big)(x) = \Big(\sum_{i = 0}^{r} c_i f^i\Big)(0) = 0,
\end{array}\]

特别地 $\ker \varphi$ 对乘以 $K[T]$ 中元素封闭.但我们知道, $K[T]$ 是主理想环, 因此存在唯一确定的次数最小的首一多项式

\[
p(T) = T^r + c_1 T^{r-1} + \ldots + c_r \in \ker \varphi,\]

并且它生成该理想; 因为 $\ker \varphi$ 并不仅由零多项式组成, 因为 $K[T]$ 作为 $K$-向量空间维数为无穷, 而 $V$ 为有限维.由此得到方程

\[
p(f)(x) = f^r(x) + c_1 f^{r-1}(x) + \ldots + c_r x = 0,\]

它以归纳的方式表明 $U$ 已经由 $x, f(x), \ldots, f^{r-1}(x)$ 生成, 并且由于 $\ker \varphi$ 中不存在次数 $< r$ 的非平凡多项式, 这些元素甚至构成 $U$ 的一个基.关于这个基, $f|_U$ 就由多项式 $p$ 的所谓\emph{友矩阵}(\textup{Begleitmatrix})\index{youjuzhen@友矩阵(Begleitmatrix)}

\[
\left(\begin{array}{cccccc}
0 & & & & & -c_0 \\
1 & 0 & & & & -c_1 \\
 & 1 & \ddots & & & -c_2 \\
 & & \ddots & \ddots & & \vdots \\
 & & & \ddots & 0 & \vdots \\
 & & & & 1 & -c_{n-1}
\end{array}\right)\]

来描述.

为了完全弄清开头的问题 2, 还要研究 $V$ 在多大程度上能表示为 $f$-循环线性子空间的直和.事实上这总是可能的, 而且可以用所谓 $f$-不可分解的(同样也是 $f$-循环的)线性子空间来实现, 这些子空间不再允许进一步分解为 $f$-\emph{不变子空间}(\textup{invarianter Untervektorraum})\index{bubianzikongjian@不变子空间(invarianter Untervektorraum)}的直和, 并且配有来自 $K[T]$ 的,由 $f$ 唯一确定的首一多项式; 这是一个深刻的结论, 其证明需要一定的功夫.为了揭示这个结论得以成立的真正原因, 我们将处理所谓\emph{初等因子理论}(\textup{Elementarteilertheorie})\index{chudengyinzililun@初等因子理论(Elementarteilertheorie)}, 而且是对主理想环上的模来处理的.环上的模的定义与域上的向量空间在形式上完全一样, 只是把纯量域换成环.例如, 应当把带自同态 $f \colon V \longrightarrow V$ 的 $K$-向量空间 $V$ 看作多项式环 $K[T]$ 上的模, 其中对 $x \in V$ 把乘积 $Tx$ 定义为 $f(x)$, 这一点正是由上面关于 $f$-循环线性子空间的考虑所提示的.在此意义上, 由向量 $x \in V$ 生成的 $f$-循环线性子空间 $U \subset V$ 应看作 $V$ 中由 $x$ 生成的 $K[T]$-\emph{子模}(\textup{Untermodul})\index{zimo@子模(Untermodul)}.

虽然模可以解释为环上的"向量空间", 但模与域上的向量空间之间仍有重大差别, 它们是由所谓扭(torsion)现象造成的.因为对环 $R$ 上的模 $M$, 一般存在非零元素 $r \in R$ 与 $m \in M$ 使 $rm = 0$, 此时 $r$ 自然不可能是单位, 否则会推出 $m = r^{-1}(rm) = 0$.特别地, 这样的模不可能有基.若对每个 $m \neq 0$(其中 $m \in M$) 都存在 $r \neq 0$(其中 $r \in R$) 使 $rm = 0$, 则称 $M$ 为\emph{扭模}(\textup{Torsionsmodul})\index{niumo@扭模(Torsionsmodul)}.例如在上述情形中, $V$ 作为 $K[T]$-模也是扭模, 因为 $V$ 有限维, 所以存在非平凡的多项式 $p \in K[T]$ 使 $p \cdot V = 0$.我们将特别看到, 又存在唯一确定的次数最小的首一多项式 $p_f \in K[T]$ 满足 $p_f \cdot V = 0$.称 $p_f$ 为 $f$ 的最小多项式, 并且我们将用凯莱--哈密顿(Cayley-Hamilton)定理证明, $p_f$ 总是特征多项式 $\chi_f$ 的一个因子, 从而次数 $\leq n$.

利用初等因子理论得到的 $V$ 到 $f$-循环的即 $f$-不变的子空间的分解, 我们将最终用来借助典范对应的矩阵描述 $f$.例如若 $V = \bigoplus_{i = 1}^{r} U_i$ 是 $f$-循环线性子空间的分解, 其中在 $K[T]$-模的意义下 $U_i \simeq K[T]/(p_i)$ 带有首一多项式 $p_i \in K[T]$, 则可以对每个 $U_i$ 选取一个 $K$-基, 使得 $f|_{U_i}$ 关于该基由 $p_i$ 的友矩阵 $A(p_i)$ 表示.把各个 $U_i$ 的基合成为 $V$ 的总基, 则描述 $f$ 的矩阵具有 $\operatorname{Diag}(A(p_1), \ldots, A(p_r))$ 的形式, 即一种"对角矩阵", 其对角线上排列着方格 $A(p_i)$.若从描述 $f$ 的任意矩阵 $A \in K^{n \times n}$ 出发, 并且诸多项式 $p_i$ 都是素幂, 则 $A' = \operatorname{Diag}(A(p_1), \ldots, A(p_r))$ 已经是 $A$ 的所谓\emph{一般标准形}(\textup{allgemeine Normalform})\index{yibanbiaozhunxing@一般标准形(allgemeine Normalform)}.此外我们还将考虑 $A$ 的若尔当(Jordan)\emph{标准形}(\textup{Normalform})\index{biaozhunxing@标准形(Normalform)}(只要特征多项式 $\chi_f$ 分解为线性因子), 并且最后说明如何通过确定矩阵 $TE - A \in K[T]^{n \times n}$ 的\emph{初等因子}(\textup{Elementarteiler})\index{chudengyinzi@初等因子(Elementarteiler)}来显式地计算标准形; 其中 $E \in K^{n \times n}$ 为单位矩阵.

\section{特征值与特征向量}\index{tezhengzhi@特征值(Eigenwert)}\index{tezhengxiangliang@特征向量(Eigenvektor)}\label{sec:6.1}

我们现在回到域 $K$ 上向量空间的理论, 首先考虑有限维 $K$-向量空间 $V$ 与 $W$ 之间的 $K$-线性映射 $f \colon V \longrightarrow W$.若 $X = (x_1, \ldots, x_n)$ 是 $V$ 的一个基, $Y = (y_1, \ldots, y_m)$ 是 $W$ 的一个基, 则 $f$ 可以由相应的矩阵 $A_{f, X, Y}$ 描述; 参见定义~\mtref{def:3.1.3}.通过适当选取 $X$ 和 $Y$, 可以使 $A_{f, X, Y}$ 具有尽可能简单的形状.例如我们在引理~\mtref{lem:3.4.7} 中看到, 存在 $V$ 的基 $X'$ 和 $W$ 的基 $Y'$, 使得

\[
A_{f, X', Y'} = \left(\begin{array}{cc} E_r & 0 \\ 0 & 0 \end{array}\right);\]

其中 $E_r \in K^{r \times r}$ 为单位矩阵, $r$ 为 $f$ 的秩.此外, 由基变换矩阵 $A_{\operatorname{id}, X', X}$ 和 $A_{\operatorname{id}, Y', Y}$, 参见定理~\mtref{thm:3.4.4}, 有关系

\[
A_{f, X', Y'} = (A_{\operatorname{id}, Y', Y})^{-1} \cdot A_{f, X, Y} \cdot A_{\operatorname{id}, X', X}\text{.}\]

利用线性映射与矩阵之间的双射对应, 我们因此也可以说: 对秩为 $r$ 的矩阵 $A \in K^{m \times n}$, 恒存在可逆矩阵 $S \in \mathrm{Gl}(m, K)$ 和 $T \in \mathrm{Gl}(n, K)$ 使

\[
S^{-1} \cdot A \cdot T = \left(\begin{array}{cc} E_r & 0 \\ 0 & 0 \end{array}\right),\]

参见定理~\mtref{thm:3.4.8}.

在下面的叙述中, 我们要研究自同态 $f \colon V \longrightarrow V$ 的相应问题.更确切地说, 要通过适当选取 $V$ 的基 $X$, 使矩阵 $A_{f, X, X}$ 具有尽可能简单的形状.译成矩阵的语言就是: 从矩阵 $A \in K^{n \times n}$ 出发, 要构造可逆矩阵 $S \in \mathrm{Gl}(n, K)$, 使得矩阵 $S^{-1} \cdot A \cdot S$ 具有尽可能简单的形状, 例如对角矩阵

\[
\left(\begin{array}{cccc}
\lambda_1 & & & 0 \\
 & \lambda_2 & & \\
 & & \ddots & \\
0 & & & \lambda_n
\end{array}\right)\]

其中 $\lambda_1, \ldots, \lambda_n \in K$ 为任意系数.这里要指出, 这样的对角形状并非总能达到.为给我们的问题准备一种方便的表述方式, 我们说:

\begin{definition}[相似矩阵]\label{def:6.1.1}\index{xiangsi@相似(ähnlich)}
两个矩阵 $A, B \in K^{n \times n}$ 称为相似的, 如果存在可逆矩阵 $S \in \mathrm{Gl}(n, K)$ 使 $B = S^{-1} \cdot A \cdot S$.
\end{definition}

立即看到, $K^{n \times n}$ 中矩阵的相似构成一个等价关系.于是 $K^{n \times n}$ 分解为相似矩阵的不相交的类.

\begin{remark}\label{rem:6.1.2}
对矩阵 $A \in K^{n \times n}$, 考虑一个 $n$ 维 $K$-向量空间 $V$ 及其一个基 $X$, 以及(唯一确定的)自同态 $f \colon V \longrightarrow V$, 满足 $A_{f, X, X} = A$; 参见定理~\mtref{thm:3.3.2}.对另一矩阵 $B \in K^{n \times n}$, 下列条件等价:
\begin{enumerate}
\item[(i)] $A$ 与 $B$ 相似;
\item[(ii)] 存在 $V$ 的基 $X'$ 使 $A_{f, X', X'} = B$.
\end{enumerate}
\end{remark}

\begin{proof}
首先设 $A$ 与 $B$ 相似, 即存在 $S \in \mathrm{Gl}(n, K)$ 使 $B = S^{-1} \cdot A \cdot S$.按注记~\mtref{rem:3.4.1}(以及随后的说明) 把矩阵 $S$ 理解为基变换矩阵, 则得到 $V$ 的一个基 $X'$, 使 $S = A_{\operatorname{id}, X', X}$, 并且由定理~\mtref{thm:3.4.4} 得到

\[
A_{f, X', X'} = (A_{\operatorname{id}, X', X})^{-1} \cdot A_{f, X, X} \cdot A_{\operatorname{id}, X', X} = S^{-1} \cdot A \cdot S = B,\]

即条件 (ii) 成立.

反之, 若条件 (ii) 成立, 则方程

\[
A_{f, X', X'} = (A_{\operatorname{id}, X', X})^{-1} \cdot A_{f, X, X} \cdot A_{\operatorname{id}, X', X}\]

表明 $A$ 与 $B$ 相似.
\end{proof}

我们现在想研究给定的矩阵 $A \in K^{n \times n}$ 何时相似于一个对角矩阵.为此引入下面的表述:

\begin{definition}[可对角化]\label{def:6.1.3}\index{keduijiaohua@可对角化(diagonalisierbar)}
矩阵 $A \in K^{n \times n}$ 称为可对角化的, 如果它相似于某个对角矩阵.

有限维 $K$-向量空间 $V$ 的自同态 $f \colon V \longrightarrow V$ 称为可对角化的, 如果关于 $V$ 的某个基 $X$(从而, 参见推论~\mtref{cor:3.4.5}, 关于 $V$ 的一切基) 的描述矩阵 $A_{f, X, X}$ 可对角化.
\end{definition}

由注记~\mtref{rem:6.1.2} 可以看出, 有限维 $K$-向量空间 $V$ 的自同态 $f \colon V \longrightarrow V$ 可对角化, 当且仅当存在 $V$ 的基 $X = (x_1, \ldots, x_n)$ 使

\[
A_{f, X, X} = \left(\begin{array}{cccc}
\lambda_1 & & & 0 \\
 & \lambda_2 & & \\
 & & \ddots & \\
0 & & & \lambda_n
\end{array}\right),\]

从而对某些常数 $\lambda_i \in K$, $i = 1, \ldots, n$, 有 $f(x_i) = \lambda_i x_i$.在这种联系下我们将使用下面的术语:

\begin{definition}[特征值与特征向量]\label{def:6.1.4}\index{tezhengzhi@特征值(Eigenwert)}\index{tezhengxiangliang@特征向量(Eigenvektor)}
设 $f \colon V \longrightarrow V$ 是 $K$-向量空间 $V$ 的自同态.常数 $\lambda \in K$ 称为 $f$ 的特征值, 如果存在向量 $a \in V - \{0\}$ 使 $f(a) = \lambda a$.此时称 $a$ 为 $f$ 的属于特征值 $\lambda$ 的特征向量.对矩阵 $A \in K^{n \times n}$, 特征值与特征向量定义为相应的线性映射 $K^n \longrightarrow K^n$, $x \longmapsto Ax$ 的特征值与特征向量.
\end{definition}

按定义, 特征向量总是不等于 $0$ 的, 于是可以表述:

\begin{remark}\label{rem:6.1.5}
有限维 $K$-向量空间 $V$ 的自同态 $f \colon V \longrightarrow V$ 可对角化, 当且仅当 $V$ 中存在由关于 $f$ 的特征向量组成的基.
\end{remark}

作为用矩阵描述线性映射(参见定理~\mtref{thm:3.1.7})的应用, 得到:

\begin{remark}\label{rem:6.1.6}
设 $f \colon V \longrightarrow V$ 是带基 $X$ 的有限维 $K$-向量空间 $V$ 的自同态.对 $\lambda \in K$, 下列条件等价:
\begin{enumerate}
\item[(i)] $\lambda$ 是 $f$ 的特征值;
\item[(ii)] $\lambda$ 是 $A_{f, X, X}$ 的特征值.
\end{enumerate}
\end{remark}

\begin{proof}
设 $\dim_K V = n$.我们利用推论~\mtref{cor:3.1.8} 中的交换图.其中 $\kappa_X$ 是把向量 $v \in V$ 送到相应的坐标列向量 $v_X \in K^n$ 的那个同构, 而 $\tilde{f} \colon K^n \longrightarrow K^n$ 是由 $u \longmapsto A_{f, X, X} \cdot u$ 定义的映射.若 $a \in V$ 是 $f$ 的属于特征值 $\lambda \in K$ 的特征向量, 则特别有 $a \neq 0$, 从而也有 $\kappa_X(a) \neq 0$.再由上面交换图的交换性得到

\begin{center}
\begin{tikzcd}[row sep=huge, column sep=huge]
V \arrow[r, "f"] \arrow[d, "\kappa_X"'] & V \arrow[d, "\kappa_X"] \\
K^n \arrow[r, "\tilde{f}"'] & K^n
\end{tikzcd}
\end{center}

\[
\tilde{f}\big(\kappa_X(a)\big) = \kappa_X\big(f(a)\big) = \kappa_X(\lambda a) = \lambda \kappa_X(a),\]

即 $\kappa_X(a)$ 是 $\tilde{f}$ 的属于特征值 $\lambda$ 的特征向量.反之, 若 $b \in K^n$ 是 $\tilde{f}$ 的属于特征值 $\lambda$ 的特征向量, 则相应地推出 $\kappa_X^{-1}(b) \in V$ 是 $f$ 的属于特征值 $\lambda$ 的特征向量.
\end{proof}

特别地, 由注记~\mtref{rem:6.1.2} 或直接的简单计算可见:

\begin{remark}[相似矩阵具有相同的特征值]\label{rem:6.1.7}
相似矩阵具有相同的特征值.
\end{remark}

作为例子, 我们证明矩阵

\[
A = \left(\begin{array}{cc} 0 & 1 \\ -1 & 0 \end{array}\right) \in \mathbb{R}^{2 \times 2}, \qquad B = \left(\begin{array}{cc} 1 & 0 \\ 1 & 1 \end{array}\right) \in \mathbb{C}^{2 \times 2}\]

不可对角化.由于方程组

\[
\alpha_2 = \lambda \alpha_1, \qquad -\alpha_1 = \lambda \alpha_2\]

对非平凡向量 $(\alpha_1, \alpha_2)^t \in \mathbb{R}^2$ 总是归结为在 $\mathbb{R}$ 中不可解的方程 $\lambda^2 = -1$, 我们看到 $A$ 在 $\mathbb{R}^{2 \times 2}$ 中不可对角化, 因为相应的线性映射 $\mathbb{R}^2 \longrightarrow \mathbb{R}^2$, $x \longmapsto Ax$ 没有特征值.然而, 若把 $A$ 看作 $\mathbb{C}^{2 \times 2}$ 中的矩阵, 图像就改变了, 因为由 $A$ 给出的 $\mathbb{C}$-线性映射 $\mathbb{C}^2 \longrightarrow \mathbb{C}^2$, $x \longmapsto Ax$ 关于基 $\binom{1}{i}$, $\binom{1}{-i}$ 由对角矩阵描述.此外, 方程组

\[
\alpha_1 = \lambda \alpha_1, \qquad \alpha_1 + \alpha_2 = \lambda \alpha_2\]

表明矩阵 $B$ 至多只有特征值 $\lambda = 1$.若 $B$ 可对角化, 则 $B$ 必相似于单位矩阵, 从而已经与单位矩阵一致, 因为单位矩阵由于平凡的原因只与自身相似.


\begin{theorem}\label{thm:6.1.8}
设 $f \colon V \longrightarrow V$ 是有限维 $K$-向量空间 $V$ 的自同态.若 $a_1, \ldots, a_r \in V$ 是属于两两不同的特征值 $\lambda_1, \ldots, \lambda_r$ 的特征向量, 则 $a_1, \ldots, a_r$ 线性无关.
\end{theorem}

\begin{proof}
我们对 $r$ 作归纳, 其中可设 $r \geq 1$.$r = 1$ 的情形是清楚的, 因为特征向量按定义总是不等于 $0$.设 $r > 1$, 并设

\[
\sum_{i = 1}^{r} \alpha_i a_i = 0\]

其中系数 $\alpha_1, \ldots, \alpha_r \in K$.于是有

\[
\sum_{i = 1}^{r} \lambda_i \alpha_i a_i = \sum_{i = 1}^{r} \alpha_i f(a_i) = f\Big(\sum_{i = 1}^{r} \alpha_i a_i\Big) = f(0) = 0,\]

但也有

\[
\sum_{i = 1}^{r} \lambda_1 \alpha_i a_i = 0\]

从而

\[
\sum_{i = 2}^{r} (\lambda_i - \lambda_1) \alpha_i a_i = 0\text{.}\]

现在 $a_2, \ldots, a_r$ 总共是 $r - 1$ 个属于两两不同特征值的特征向量, 因而按归纳假设线性无关.于是得到 $(\lambda_i - \lambda_1)\alpha_i = 0$, 从而对 $i = 2, \ldots, r$ 有 $\alpha_i = 0$.方程 $\sum_{i = 1}^{r} \alpha_i a_i = 0$ 于是表明项 $\alpha_1 a_1$ 也消失, 并且由于 $a_1 \neq 0$ 甚至系数 $\alpha_1$ 也消失.因此 $a_1, \ldots, a_r$ 如所断言线性无关.
\end{proof}

\begin{corollary}\label{cor:6.1.9}
有限维 $K$-向量空间 $V$ 的自同态 $f \colon V \longrightarrow V$ 至多有 $n = \dim_K V$ 个不同的特征值.
\end{corollary}

\begin{corollary}\label{cor:6.1.10}
若有限维 $K$-向量空间 $V$ 的自同态 $f \colon V \longrightarrow V$ 恰好有 $n = \dim_K V$ 个不同的特征值, 则 $f$ 可对角化.
\end{corollary}

\begin{proof}
设 $\lambda_1, \ldots, \lambda_n \in K$ 是 $f$ 的两两不同的特征值, 并设 $a_1, \ldots, a_n \in V$ 是相应的特征向量.按定理~\mtref{thm:6.1.8} 它们线性无关, 因而由于 $n = \dim_K V$ 构成 $V$ 的一个基 $X$.相应的矩阵 $A_{f, X, X}$ 是以 $\lambda_1, \ldots, \lambda_n$ 为对角元素的对角矩阵.
\end{proof}

我们想把定理~\mtref{thm:6.1.8} 再稍加推广.

\begin{definition}[特征空间]\label{def:6.1.11}\index{tezhengkongjian@特征空间(Eigenraum)}
设 $\lambda \in K$ 是向量空间自同态 $f \colon V \longrightarrow V$ 的特征值.则称

\[
V_\lambda := \ker(f - \lambda \operatorname{id}) = \{a \in V;\; f(a) = \lambda a\}\]

为 $f$ 的属于特征值 $\lambda$ 的特征空间.
\end{definition}

\begin{corollary}\label{cor:6.1.12}
设 $f \colon V \longrightarrow V$ 是有限维 $K$-向量空间 $V$ 的自同态, 并设 $\lambda_1, \ldots, \lambda_r$ 是 $f$ 的全部(两两不同的)特征值.若 $V' = \sum_{i = 1}^{r} V_{\lambda_i}$ 是由相应的特征空间生成的 $V$ 的子空间, 则

\[
V' = \bigoplus_{i = 1}^{r} V_{\lambda_i};\]

即这个和是直和.此外, $f$ 可对角化, 当且仅当 $V = V'$, 即当且仅当 $V$ 由 $f$ 的特征空间生成.
\end{corollary}

\begin{proof}
由于属于两两不同特征值的特征向量线性无关, 和 $\sum_{i = 1}^{r} v_i$(其中 $v_i \in V_{\lambda_i}$) 只在所有 $v_i$ 都消失时才消失.这意味着 $V'$ 是特征空间 $V_{\lambda_i}$ 的直和.若 $f$ 可对角化, 则 $V$ 有一个由特征向量组成的基, 于是 $V = V'$.反之, 若 $V$ 可以表示为特征空间 $V_{\lambda_i}$ 的直和, 则在这些特征空间的每一个中选取一个基.所有这些元素的全体构成 $V$ 的一个基, 它完全由 $f$ 的特征向量组成, 即 $f$ 可对角化.
\end{proof}

\begin{corollary}\label{cor:6.1.13}
对对角矩阵

\[
A = \left(\begin{array}{cccc}
\lambda_1 & & & 0 \\
 & \lambda_2 & & \\
 & & \ddots & \\
0 & & & \lambda_n
\end{array}\right),\]

$\lambda_1, \ldots, \lambda_n$ 是 $A$ 仅有的特征值.
\end{corollary}

\begin{proof}
我们考虑线性映射 $f \colon K^n \longrightarrow K^n$, $x \longmapsto Ax$.显然 $\lambda_1, \ldots, \lambda_n$ 是 $A$ 即 $f$ 的特征值.为避免重复, 把 $\{\lambda_1, \ldots, \lambda_n\}$ 写成 $\{\lambda_1', \ldots, \lambda_s'\}$, 其中 $\lambda_1', \ldots, \lambda_s'$ 两两不同.$A$ 的对角形状表明 $K^n$ 中存在一个基, 即典范基 $e_1, \ldots, e_n$, 使得 $e_i$ 分别是 $f$ 的属于特征值 $\lambda_i$ 的特征向量.对 $j = 1, \ldots, s$, 设 $U_{\lambda_j'} \subset V$ 是由所有满足 $\lambda_i = \lambda_j'$ 的指标 $i$ 对应的 $e_i$ 生成的线性子空间.则有分解

\[
V = \bigoplus_{j = 1}^{s} U_{\lambda_j'}, \qquad U_{\lambda_j'} \subset V_{\lambda_j'},\]

其中 $V_{\lambda_j'}$ 分别是 $f$ 的属于特征值 $\lambda_j'$ 的特征空间.与推论~\mtref{cor:6.1.12} 中的分解比较, 得到对 $j = 1, \ldots, s$ 有 $U_{\lambda_j'} = V_{\lambda_j'}$, 并且还表明除了 $\lambda_1', \ldots, \lambda_s'$ 之外 $A$ 不可能再有其他特征值.
\end{proof}

\aufgaben

下面 $V$ 恒为域 $K$ 上的有限维向量空间.

\noindent\textbf{习题 1.} 设 $a, b \in V$ 是自同态 $f \colon V \longrightarrow V$ 的特征向量.研究在哪些情形下 $a - b$ 也是 $f$ 的特征向量.

\medskip
\noindent\textbf{习题 2.} 设 $\lambda \in K$ 是自同态 $f \colon V \longrightarrow V$ 的特征值.证明: 对多项式 $q \in K[T]$, $q(\lambda)$ 分别是 $q(f)$ 的特征值.

\medskip
\noindent\textbf{习题 3.} 对矩阵

\[
A = \left(\begin{array}{cccc} 2 & 1 & 0 & 1 \\ 0 & 2 & 0 & 1 \\ 0 & 0 & 2 & 1 \\ 0 & 0 & 0 & 2 \end{array}\right) \in \mathbb{R}^{4 \times 4}\]

计算全部特征值与相应的特征空间.$A$ 可对角化吗? (AT 347)

\medskip
\noindent\textbf{习题 4.} 设矩阵 $A, B \in K^{n \times n}$ 相似.直接证明:
\begin{enumerate}
\item[(i)] 元素 $\lambda \in K$ 是 $A$ 的特征值, 当且仅当它是 $B$ 的特征值;
\item[(ii)] 对 $A$ 即 $B$ 的特征值 $\lambda \in K$, 有 $\dim_K V_{A, \lambda} = \dim_K V_{B, \lambda}$, 其中 $V_{A, \lambda}$ 表示线性映射 $K^n \longrightarrow K^n$, $x \longmapsto Ax$ 的属于 $\lambda$ 的特征空间; 对 $V_{B, \lambda}$ 类似.
\end{enumerate}

\medskip
\noindent\textbf{习题 5.} 设 $A, B \in K^{n \times n}$ 相似.则对多项式 $q \in K[T]$, 矩阵 $q(A)$ 与 $q(B)$ 也相似.

\medskip
\noindent\textbf{习题 6.} 两个自同态 $f, g \colon V \longrightarrow V$ 称为相似的, 如果存在自同构 $h \colon V \longrightarrow V$ 使 $g = h^{-1} \circ f \circ h$.证明: $f$ 与 $g$ 相似, 当且仅当对给定基 $X$ 描述矩阵 $A_{f, X, X}$ 与 $A_{g, X, X}$ 相似.

\medskip
\noindent\textbf{习题 7.} 对 $K$-向量空间之间的线性映射的交换图

\begin{center}
\begin{tikzcd}[row sep=huge, column sep=huge]
V \arrow[r, "f"] \arrow[d, "h"'] & V \arrow[d, "h"] \\
W \arrow[r, "g"'] & W
\end{tikzcd}
\end{center}

证明:
\begin{enumerate}
\item[(i)] 若 $h$ 是单射, 则 $f$ 的每个特征值也是 $g$ 的特征值;
\item[(ii)] 若 $h$ 是满射, 则 $g$ 的每个特征值也是 $f$ 的特征值.
\end{enumerate}

构造简单的例子, 表明在上述断言中"单射"与"满射"的假设不可缺少.(AT 349)

\section{最小多项式与特征多项式}\index{zuixiaoduoxiangshi@最小多项式(Minimalpolynom)}\index{tezhengduoxiangshi@特征多项式(charakteristisches Polynom)}\label{sec:6.2}

正如上一节所见, 自同态或矩阵的可对角化问题与确定相应的特征值和特征向量问题密切相关.因此我们首先处理特征值的计算.在本节中, 总设 $V$ 是有限维数 $n$ 的 $K$-向量空间.

\begin{theorem}\label{thm:6.2.1}
设 $f \colon V \longrightarrow V$ 是自同态.对 $\lambda \in K$, 下列条件等价:
\begin{enumerate}
\item[(i)] $\lambda$ 是 $f$ 的特征值;
\item[(ii)] $\ker(\lambda \operatorname{id} - f) \neq 0$;
\item[(iii)] $\lambda \operatorname{id} - f$ 不可逆;
\item[(iv)] $\det(\lambda \operatorname{id} - f) = 0$.
\end{enumerate}
\end{theorem}

\begin{proof}
设 $\lambda$ 是 $f$ 的特征值.则存在属于 $\lambda$ 的特征向量, 即存在 $a \in V - \{0\}$ 使 $f(a) = \lambda a$.由此得到 $a \in \ker(\lambda \operatorname{id} - f)$, 从而特别有 $\ker(\lambda \operatorname{id} - f) \neq 0$.反之, $\ker(\lambda \operatorname{id} - f)$ 中每个不等于 $0$ 的向量都是属于特征值 $\lambda$ 的特征向量.条件 (i) 与 (ii) 因此等价.

进一步, 等价 (ii) $\longleftrightarrow$ (iii) 由推论~\mtref{cor:2.1.11} 得出, 等价 (iii) $\longleftrightarrow$ (iv) 由定理~\mtref{thm:4.3.3}(iv) 得出.
\end{proof}

作为例子, 我们想确定矩阵

\[
A = \left(\begin{array}{cc} 1 & 4 \\ 1 & 1 \end{array}\right) \in \mathbb{R}^{2 \times 2}\]

的全部特征值.若 $E$ 表示 $\mathbb{R}^{2 \times 2}$ 中的单位矩阵, 则对 $\lambda \in \mathbb{R}$ 有

\[
\det(\lambda E - A) = \det\left(\begin{array}{cc} \lambda - 1 & -4 \\ -1 & \lambda - 1 \end{array}\right) = \lambda^2 - 2\lambda - 3\text{.}\]

方程 $\det(\lambda E - A) = 0$ 因此等价于 $\lambda = 3$ 或 $\lambda = -1$.故 $3, -1$ 是 $A$ 的特征值, 并且由推论~\mtref{cor:6.1.10} 可见 $A$ 可对角化.事实上 $A$ 相似于矩阵

\[
\left(\begin{array}{cc} 3 & 0 \\ 0 & -1 \end{array}\right) \in \mathbb{R}^{2 \times 2}\text{.}\]

现在可以通过解线性方程组

\[
(3E - A)x = 0 \quad \text{以及} \quad (-E - A)x = 0\]

容易确定属于特征值 $3, -1$ 的特征空间.

我们想对矩阵 $A = (\alpha_{ij})_{ij} \in K^{n \times n}$ 和常数 $\lambda \in K$ 更精确地计算 $\det(\lambda E - A)$; $E$ 现在是 $K^{n \times n}$ 中的单位矩阵.如引理~\mtref{lem:4.2.4} 中定义的那样,

\[
\begin{array}{l}
\det(\lambda E - A) = \det\big((\lambda \delta_{ij} - \alpha_{ij})_{ij}\big) \\
\qquad = \sum_{\pi \in \mathfrak{S}_n} \sgn \pi \cdot \prod_{i = 1}^{n} (\lambda \delta_{\pi(i), i} - \alpha_{\pi(i), i}),
\end{array}\]

于是可以作下面的定义:

\begin{definition}[矩阵的特征多项式]\label{def:6.2.2}\index{tezhengduoxiangshi@特征多项式(charakteristisches Polynom)}
设 $A = (\alpha_{ij})_{ij} \in K^{n \times n}$.则称

\[
\chi_A = \sum_{\pi \in \mathfrak{S}_n} \sgn \pi \cdot \prod_{i = 1}^{n} (T \delta_{\pi(i), i} - \alpha_{\pi(i), i}) \in K[T]\]

为 $A$ 的特征多项式.
\end{definition}

特别有 $\chi_A(\lambda) = \det(\lambda E - A)$ 对一切 $\lambda \in K$ 成立, 并且 $\chi_A$ 在 $K$ 中的零点正好是 $A$ 的特征值.

\begin{theorem}\label{thm:6.2.3}
矩阵 $A \in K^{n \times n}$ 的特征多项式 $\chi_A \in K[T]$ 是次数为 $n$ 的首一多项式.若

\[
\chi_A = \sum_{i = 0}^{n} c_{n-i} T^i, \qquad c_i \in K,\]

则 $c_0 = 1$, $c_1 = -\operatorname{Spur} A = -\sum_{i = 1}^{n} \alpha_{ii}$, 且 $(-1)^n c_n = \det(A)$, 其中对角元素 $\alpha_{ii}$ 的和称为 $A$ 的\emph{迹}(\textup{Spur})\index{ji@迹(Spur)}.
\end{theorem}

\begin{proof}
作为 $T$ 中线性多项式的 $n$ 重乘积之和, $\chi_A$ 的次数 $\leq n$, 并且

\[
c_n = \chi_A(0) = \det(-A) = (-1)^n \det(A),\]

即 $(-1)^n c_n = \det(A)$.此外, 分解

\[
\chi_A = \prod_{i = 1}^{n} (T - \alpha_{ii}) + \sum_{\substack{\pi \in \mathfrak{S}_n \\ \pi \neq \operatorname{id}}} \sgn \pi \cdot \prod_{i = 1}^{n} (T \delta_{\pi(i), i} - \alpha_{\pi(i), i})\]

中的第二项次数 $\leq n - 2$, 因为只对 $\pi \neq \operatorname{id}$ 求和.对 $\pi \neq \operatorname{id}$, 至少存在两个不同指标 $i, j \in \{1, \ldots, n\}$ 使 $\pi(i) \neq i$, $\pi(j) \neq j$, 因而对这些指标表达式 $T \delta_{\pi(i), i}$, $T \delta_{\pi(j), j}$ 消失.因此 $\chi_A$ 中次数为 $n$ 与 $n - 1$ 的系数, 即 $c_0$ 与 $c_1$, 与 $\prod_{i = 1}^{n}(T - \alpha_{ii})$ 中次数为 $n$ 与 $n - 1$ 的系数一致, 于是如所断言 $c_0 = 1$ 以及 $c_1 = -\sum_{i = 1}^{n} \alpha_{ii}$.
\end{proof}

矩阵 $A \in K^{n \times n}$ 的特征多项式 $\chi_A \in K[T]$ 也可以通过方程 $\chi_A = \det(TE - A)$ 来定义, 不过此时要取系数在多项式环 $K[T]$ 中的矩阵的行列式.由于我们迄今只考虑系数在一个域中的矩阵的行列式, 并且只对这种情形证明过行列式的运算规则, 我们使用一个技巧.类似于由整数构造有理数域 $\mathbb{Q}$ 那样, 由 $K[T]$ 构造所谓\emph{有理函数域}(\textup{rationaler Funktionenkörper})\index{youlihanshuyu@有理函数域(rationaler Funktionenkörper)} $K(T)$, 它由 $K[T]$ 中多项式的全部商组成.则 $K[T]$ 是域 $K(T)$ 的子环, 并且可以把 $TE - A$ 看作 $K(T)^{n \times n}$ 中的矩阵.特别地, $\det(TE - A)$ 是良定义的, 并且计算这个行列式时允许使用已知的展开定理或其他行列式运算规则.但也要指出, 行列式理论也可以在任意交换环上展开, 不过我们在此不再深入讨论.

\begin{theorem}\label{thm:6.2.4}
若 $A, B \in K^{n \times n}$ 相似, 则 $\chi_A = \chi_B$.
\end{theorem}

\begin{proof}
设 $S \in \mathrm{Gl}(n, K)$ 满足 $B = S^{-1} A S$.则利用行列式的乘性,

\[
\begin{array}{r l}
& \chi_B = \det(TE - S^{-1} A S) = \det\bigl(S^{-1}(TE - A)S\bigr) \\
& \qquad = \det(S^{-1}) \cdot \det(TE - A) \cdot \det(S) = \det(TE - A) \\
& \qquad = \chi_A\text{.}
\end{array}\]

\end{proof}

\begin{corollary}[相似矩阵具有相同的迹]\label{cor:6.2.5}\index{ji@迹(Spur)}
相似矩阵具有相同的迹.
\end{corollary}

由于相似矩阵的特征多项式一致, 现在可以利用注记~\mtref{rem:6.1.2} 定义自同态 $f \colon V \longrightarrow V$ 的特征多项式.

\begin{definition}[自同态的特征多项式与迹]\label{def:6.2.6}\index{ji@迹(Spur)}
设 $f \colon V \longrightarrow V$ 是自同态, $X$ 是有限维 $K$-向量空间 $V$ 的一个基.则称

\[
\chi_f = \chi_{A_{f, X, X}} \in K[T]\]

为 $f$ 的特征多项式, 并称

\[
\operatorname{Spur} f = \operatorname{Spur} A_{f, X, X}\]

为 $f$ 的迹.对平凡情形 $V = 0$, 有 $\chi_f = 1$ 和 $\operatorname{Spur} f = 0$.
\end{definition}

由注记~\mtref{rem:6.1.2} 和定理~\mtref{thm:6.2.4} 推出, $\chi_f$ 与 $\operatorname{Spur} f$ 不依赖于 $X$ 的特殊选取.此外, 由定理~\mtref{thm:6.2.1} 我们可以确定:

\begin{theorem}\label{thm:6.2.7}
元素 $\lambda \in K$ 是自同态 $f \colon V \longrightarrow V$ 的特征值, 当且仅当 $\lambda$ 是特征多项式 $\chi_f$ 的零点.
\end{theorem}

接下来我们要看出, 自同态的可对角化性可以由相应特征多项式的某些性质读出.

\begin{theorem}\label{thm:6.2.8}
对自同态 $f \colon V \longrightarrow V$, 下列条件等价:
\begin{enumerate}
\item[(i)] $f$ 可对角化;
\item[(ii)] $\chi_f$ 完全分解为线性因子, 例如

\[
\chi_f = \prod_{i = 1}^{r} (T - \lambda_i)^{n_i},\]

并且对属于特征值 $\lambda_i$ 的特征空间 $V_{\lambda_i}$ 有

\[
\dim V_{\lambda_i} = n_i\text{.}\]
\end{enumerate}
\end{theorem}

\begin{proof}
设 $\dim_K V = n$.我们从蕴含 (i) $\Longrightarrow$ (ii) 开始, 并假设 $f$ 可对角化.则存在由 $f$ 的特征向量组成的 $V$ 的基 $X$, 即

\[
A_{f, X, X} = \left(\begin{array}{cccc}
\lambda_1 & & & 0 \\
 & \lambda_2 & & \\
 & & \ddots & \\
0 & & & \lambda_n
\end{array}\right)\text{.}\]

特别推出 $\chi_f = \prod_{i = 1}^{n}(T - \lambda_i)$.把相同的因子合并为幂, 写成 $\chi_f = \prod_{i = 1}^{r}(T - \lambda_i)^{n_i}$, 则得到 $\dim_K(V_{\lambda_i}) \geq n_i$.由于 $V = \bigoplus_{i = 1}^{r} V_{\lambda_i}$(参见推论~\mtref{cor:6.1.12} 与推论~\mtref{cor:6.1.13}), 于是

\[
n = \sum_{i = 1}^{r} \dim_K(V_{\lambda_i}) \geq \sum_{i = 1}^{r} n_i = n\]

从而对一切 $i = 1, \ldots, r$ 有 $\dim_K(V_{\lambda_i}) = n_i$.

反之, 若条件 (ii) 成立, 则

\[
\sum_{i = 1}^{r} \dim_K(V_{\lambda_i}) = \sum_{i = 1}^{r} n_i = n\text{.}\]

按推论~\mtref{cor:6.1.12}, 对由特征空间 $V_{\lambda_i}$ 生成的子空间 $V' \subset V$ 有 $V' = \bigoplus_{i = 1}^{r} V_{\lambda_i}$, 即 $\dim_K V' = n = \dim_K V$, 从而 $V' = V$.再由推论~\mtref{cor:6.1.12}, $f$ 可对角化.
\end{proof}

作为定理~\mtref{thm:6.2.8} 的应用例子, 我们考虑一个自同态 $f \colon V \longrightarrow V$, 它关于某个合适的基由形状为

\[
A = \left(\begin{array}{cccc}
\lambda & & & * \\
 & \lambda & & \\
 & & \ddots & \\
0 & & & \lambda
\end{array}\right) \in K^{n \times n}\]

的\emph{三角矩阵}(\textup{Dreiecksmatrix})\index{sanjiaojuzhen@三角矩阵(Dreiecksmatrix)}描述.则 $\chi_f = \chi_A = (T - \lambda)^n$, 并且 $\lambda$ 是 $f$ 即 $A$ 的唯一特征值.若 $A$ 不是对角矩阵, 则 $\lambda \operatorname{id} - f$ 不是零映射, 从而特征空间 $V_\lambda = \ker(\lambda \operatorname{id} - f)$ 真含于 $V$.按定理~\mtref{thm:6.2.8}, $f$ 即 $A$ 因此不可对角化.但我们也可以直接看出这一点: 若 $A$ 可对角化, 则 $A$ 相似于 $\lambda E$, 其中 $E \in K^{n \times n}$ 为单位矩阵.由于 $E$ 从而 $\lambda E$ 与 $K^{n \times n}$ 中一切矩阵可交换, $\lambda E$ 只能与自身相似, 因而必已有 $A = \lambda E$.


除自同态 $f$ 的特征多项式 $\chi_f$ 之外, 还可以考虑与 $f$ 相关的另一个多项式, 即所谓最小多项式.为定义它, 把自同态环 $\operatorname{End}_K(V)$ 在环同态

\[
K \longrightarrow \operatorname{End}_K(V), \qquad c \longmapsto c \cdot \operatorname{id}_V\]

下看作 $K$-代数, 并利用下面的结果:

\begin{theorem}\label{thm:6.2.9}
对有限维 $K$-向量空间 $V$ 的自同态 $f \colon V \longrightarrow V$, 构造把 $f$ 代入变量 $T$ 的 $K$-代数同态; 参见定理~\mtref{thm:5.1.7}:

\[
\varphi_f \colon K[T] \longrightarrow \operatorname{End}_K(V), \quad p \longmapsto p(f)\text{.}\]

则 $\ker \varphi_f$ 不等于零, 并且存在首一多项式 $p_f \in K[T]$ 使 $\ker \varphi_f = (p_f)$.$p_f$ 是 $K[T]$ 中满足 $p_f(f) = 0$ 的唯一确定的次数最小的首一多项式.
\end{theorem}

\begin{proof}
忘掉环乘法, 可以把 $\varphi_f$ 看作 $K$-向量空间之间的同态.于是 $\dim_K(K[T]) = \infty$, 而对 $n = \dim_K(V)$ 按定理~\mtref{thm:3.3.2} 有

\[
\dim_K\big(\operatorname{End}_K(V)\big) = n^2 < \infty\text{.}\]

后者导致 $\ker \varphi_f \neq 0$, 即 $\varphi_f$ 不可能单射.更精确地说, 元素 $\varphi_f(T^0), \ldots, \varphi_f(T^{n^2})$ 出于维数的原因在 $\operatorname{End}_K(V)$ 中线性相关, 因此 $\ker \varphi_f$ 含有一个次数 $\leq n^2$ 的非平凡多项式.但 $\ker \varphi_f$ 是 $K[T]$ 中的理想, 而 $K[T]$ 是主理想环; 参见推论~\mtref{cor:5.2.7}.因此存在非平凡多项式 $p_f \in K[T]$ 使 $\ker \varphi_f = p_f \cdot K[T]$.作为整环中主理想的生成元, $p_f$ 按注记~\mtref{rem:5.2.9} 到差一个单位为止唯一确定.但由于 $K[T]^* = K^*$, 若把这个多项式规定为首一的, 则 $p_f$ 唯一.自然 $p_f$ 就是消去 $f$ 的次数最小的首一多项式.
\end{proof}

\begin{definition}[最小多项式]\label{def:6.2.10}\index{zuixiaoduoxiangshi@最小多项式(Minimalpolynom)}
若 $f \colon V \longrightarrow V$ 是自同态, 则称按定理~\mtref{thm:6.2.9} 在 $K[T]$ 中唯一存在的,消去 $f$ 的次数最小的首一多项式为 $f$ 的最小多项式, 记为 $p_f$.相应地, 矩阵 $A \in K^{n \times n}$ 的最小多项式 $p_A$ 定义为 $K[T]$ 中消去 $A$ 的次数最小的首一多项式.
\end{definition}

在定理~\mtref{thm:6.2.9} 的证明中我们看到, 同态 $\varphi_f \colon K[T] \longrightarrow \operatorname{End}_K(V)$ 的核含有次数 $\leq n^2$ 的非平凡多项式.作为推论得到 $\operatorname{grad} p_f \leq n^2$.但这个估计还可以大大改进.

\begin{theorem}[凯莱--哈密顿]\label{thm:6.2.11}\index{kailaihamidundingli@凯莱--哈密顿定理(Cayley-Hamilton, Satz von)}
自同态 $f \colon V \longrightarrow V$ 的最小多项式 $p_f$ 总是特征多项式 $\chi_f$ 的一个因子.特别地, $\chi_f(f) = 0$, 且 $\operatorname{grad} p_f \leq \operatorname{grad} \chi_f = n$.
\end{theorem}

\begin{proof}
在定理~\mtref{thm:6.2.9} 的情形中, 只需证明 $\chi_f \in \ker \varphi_f = (p_f)$, 即 $\chi_f(f) = 0$.把它翻译为矩阵问题, 只需对矩阵 $A \in K^{n \times n}$ 证明 $\chi_A(A) = 0$.为便于计算, 再次考虑有理\emph{函数域}(\textup{Funktionenkörper})\index{hanshuyu@函数域(Funktionenkörper)} $K(T)$, 其元素是 $K[T]$ 中多项式的商.于是可以考虑 $K(T)^{n \times n}$ 的子环 $K[T]^{n \times n}$, 它由系数在 $K[T]$ 中的一切 $(n \times n)$ 矩阵组成.同态 $K(T) \longrightarrow K(T)^{n \times n}$, 把元素 $q \in K(T)$ 映到单位矩阵 $E \in K(T)^{n \times n}$ 的 $q$ 倍, 限制为环同态

\[
K[T] \longrightarrow K[T]^{n \times n}, \quad f \longmapsto f \cdot E,\]

并在 $K[T]^{n \times n}$ 上定义了 $K[T]$-代数的结构.把变量 $T \in K[T]$ 与其像 $T \cdot E \in K[T]^{n \times n}$ 等同, 则 $K[T]^{n \times n}$ 的每个元素 $\bar{M}$ 都可以写成 $M = \sum_{i \in \mathbb{N}} M_i T^i$ 的形式, 其中系数 $M_i \in K^{n \times n}$ 由 $M$ 唯一确定, 并且对几乎一切 $i \in \mathbb{N}$ 当然为零.对 $M = (\sum_{i \in \mathbb{N}} m_{\mu \nu i} T^i)_{\mu, \nu = 1, \ldots, n}$ 且系数 $m_{\mu \nu i} \in K$, 只需令 $M_i = (m_{\mu \nu i})_{\mu, \nu = 1, \ldots, n}$.这样我们可以把 $K[T]^{n \times n}$ 看作"多项式环" $K^{n \times n}[T]$, 不过要考虑到对 $n > 1$, 环 $K^{n \times n}$ 不是交换的.但如 \ref{sec:5.1} 节那样, 多项式环 $R[T]$ 也可以对非交换环 $R$ 定义, 此时必须明白: 类型为

\[
R[T] \longrightarrow R, \quad \sum a_i T^i \longmapsto \sum a_i t^i\]

的代入映射对元素 $t \in R$ 一般不再构成同态.因为 $T$ 在 $R[T]$ 中与 $R$ 中所有元素可交换, 但 $t \in R$ 未必与 $R$ 中所有其他元素可交换.

作了这些准备之后, 考虑矩阵 $A \in K^{n \times n}$ 及其特征多项式 $\chi_A \in K[T]$, 为这个证明的方便把它写成 $\chi_A = \sum_{i \in \mathbb{N}} c_i T^i$.此外, 把 $TE - A$ 看作 $K[T]^{n \times n} \subset K(T)^{n \times n}$ 中的矩阵, 并作其伴随矩阵 $(TE - A)^{\mathrm{ad}} \in K(T)^{n \times n}$; 参见定义~\mtref{def:4.4.2}.由克拉默法则~\mtref{reg:4.4.3} 得到方程

\[
(TE - A)^{\mathrm{ad}} \cdot (TE - A) = \det(TE - A) \cdot E = \chi_A(T) \cdot E\text{.}\]

这首先是 $K(T)^{n \times n}$ 中的方程, 但它也在 $K[T]^{n \times n}$ 中成立, 因为按伴随矩阵的定义, 与 $TE - A$ 一起 $(TE - A)^{\mathrm{ad}}$ 也属于 $K[T]^{n \times n}$.若写成 $(TE - A)^{\mathrm{ad}} = \sum_{i \in \mathbb{N}} A_i T^i$ 且矩阵 $A_i \in K^{n \times n}$, 则得到

\[
\begin{array}{r l}
\Big(\sum_{i \in \mathbb{N}} A_i T^i\Big) \cdot (TE - A) & = -A_0 T^0 A + \sum_{i \in \mathbb{N}} (A_i T^i T E - A_{i + 1} T^{i + 1} A) \\
& = -A_0 A + \sum_{i \in \mathbb{N}} (A_i - A_{i + 1} A) T^{i + 1} \\
& = \sum_{i \in \mathbb{N}} c_i T^i \cdot E = \chi_A(T) \cdot E,
\end{array}\]

其中

\[
-A_0 A = c_0 E, \qquad A_i - A_{i + 1} A = c_{i + 1} E \qquad \text{对 } i \in \mathbb{N}\text{.}\]

同样的计算也可以对矩阵 $B \in K^{n \times n}$ 代替变量 $T$ 进行, 只要 $A$ 与 $B$ 可交换.特别地, 在上面的方程中可以把 $T$ 换成 $B := A$, 于是如所希望得到 $\chi_A(A) = 0$, 因为项 $(TE - A)$ 在把 $T$ 换成 $A$ 后消失.
\end{proof}

我们还想考虑两个简单的例子.对单位矩阵 $E \in K^{n \times n}$, $n > 0$, 有 $\chi_E = (T - 1)^n$, $p_E = T - 1$; 对零矩阵 $0 \in K^{n \times n}$ 有 $\chi_0 = T^n$, $p_0 = T$.特别看到, 最小多项式一般与特征多项式不一致.

\aufgaben

下面 $V$ 恒为域 $K$ 上的有限维数 $n$ 的向量空间.

\noindent\textbf{习题 1.} 确定下列矩阵的特征值与相应的特征空间:

\[
A = \left(\begin{array}{cccc} 2 & 0 & 0 & 0 \\ -2 & 2 & 0 & 2 \\ 1 & 0 & 2 & 0 \\ 2 & -1 & 0 & -1 \end{array}\right)\]

$A$ 可对角化吗?

\medskip
\noindent\textbf{习题 2.} 在下列情形中确定自同态 $f \colon V \longrightarrow V$ 的最小多项式 $p_f$(AT 350):
\begin{enumerate}
\item[(i)] $V = 0$;
\item[(ii)] $f = \operatorname{id}$;
\item[(iii)] $f = 0$;
\item[(iv)] 存在线性子空间 $V_1, V_2 \subset V$ 使 $V = V_1 \oplus V_2$, 并且对 $v_i \in V_i$, $i = 1, 2$, 有 $f(v_1 + v_2) = v_1$.
\end{enumerate}

\medskip
\noindent\textbf{习题 3.} 设 $U_1, U_2 \subset V$ 是线性子空间, $f \colon V \longrightarrow V$ 是限制为自同态 $f_i \colon U_i \longrightarrow U_i$, $i = 1, 2$, 的自同态.证明:
\begin{enumerate}
\item[(i)] 若 $V = U_1 + U_2$, 则对 $f, f_1, f_2$ 的最小多项式有 $p_f = \mathrm{kgV}(p_{f_1}, p_{f_2})$;
\item[(ii)] $f$ 限制为自同态 $f_{12} \colon U_1 \cap U_2 \longrightarrow U_1 \cap U_2$, 并且对 $f_{12}, f_1, f_2$ 的最小多项式有 $p_{f_{12}} \mid \mathrm{ggT}(p_{f_1}, p_{f_2})$.一般地是否也成立等式 $p_{f_{12}} = \mathrm{ggT}(p_{f_1}, p_{f_2})$?
\end{enumerate}

\medskip
\noindent\textbf{习题 4.} 设 $K$ 代数闭.证明: 自同态 $f \colon V \longrightarrow V$ 是\emph{幂零的}(\textup{nilpotent})\index{milingde@幂零的(nilpotent)}(即满足形如 $f^r = 0$ 的方程), 当且仅当 $f$ 除 $0$ 外没有其他特征值.

\medskip
\noindent\textbf{习题 5.} 设 $f \colon V \longrightarrow V$ 是自同构.证明: 存在多项式 $q \in K[T]$ 使 $f^{-1} = q(f)$.(AT 351)

\medskip
\noindent\textbf{习题 6.} 斐波那契(Fibonacci)数列 $c_1, c_2, \ldots \in \mathbb{N}$ 由 $c_1 = c_2 = 1$ 和 $c_{n + 2} = c_{n + 1} + c_n$(对 $n \in \mathbb{N}$) 定义.给出 $c_n$ 的一个只依赖于 $n$ 的闭式表达式.(提示: 确定矩阵 $A \in \mathbb{R}^{2 \times 2}$ 使 $A \cdot (c_{n + 1}, c_n)^t = (c_{n + 2}, c_{n + 1})^t$ 对 $n \geq 1$ 成立, 以及基变换矩阵 $S \in \mathrm{Gl}(2, \mathbb{R})$, 使得 $S^{-1} \cdot A \cdot S$ 具有对角形状.)

\section{初等因子定理}\index{chudengyinzidingli@初等因子定理(Elementarteilersatz)}\index{chudengyinzi@初等因子(Elementarteiler)}\label{sec:6.3}

设 $V$ 是域 $K$ 上的向量空间, $f \colon V \longrightarrow V$ 是自同态.则在 $V$ 上定义的纯量乘法 $K \times V \longrightarrow V$, $(\alpha, v) \longmapsto \alpha \cdot v$ 延拓为外乘法

\[
K[T] \times V \longrightarrow V, \quad (p, v) \longmapsto p \cdot v := p(f)(v)\text{.}\]

其中对 $p \in K[T]$, 如通常那样 $p(f)$ 是把变量 $T$ 换成 $f$ 而由 $p$ 得到的 $\operatorname{End}_K(V)$ 中的表达式.此外, $p(f)(v)$ 是 $v$ 在自同态 $p(f) \colon V \longrightarrow V$ 下的像.容易验证, $V$ 作为加法阿贝尔群连同外乘法 $K[T] \times V \longrightarrow V$ 满足定义~\mtref{def:1.4.1} 中列出的向量空间公理, 只要不管 $K[T]$ 只是环而不是域这一点; 我们说 $V$ 是 $K[T]$-模.(注意: 与通常的用法不同, 数学中说"der Modul"或"die Moduln", 重音落在第一个音节上.)

若想研究自同态 $f \colon V \longrightarrow V$ 的标准形, 这意味着要分析上面说明的 $V$ 作为 $K[T]$-模的结构.因此我们首先整理一些关于模的基础知识.涉及的环总设为带单位元 $1$ 的交换环.

\begin{definition}[模]\label{def:6.3.1}\index{mo@模(Modul)}
设 $R$ 是带单位元 $1$ 的交换环.$R$-模是一个集合 $M$, 连同内运算 $M \times M \longrightarrow M$, $(a, b) \longmapsto a + b$, 称为加法, 以及外运算 $R \times M \longrightarrow M$, $(\alpha, a) \longmapsto \alpha \cdot a$, 称为纯量乘法, 使得:
\begin{enumerate}
\item[(i)] $M$ 关于加法是阿贝尔群;
\item[(ii)] 对一切 $\alpha, \beta \in R$, $a, b \in M$, 有 $(\alpha + \beta) \cdot a = \alpha \cdot a + \beta \cdot a$, $\alpha \cdot (a + b) = \alpha \cdot a + \alpha \cdot b$, 即加法与纯量乘法是分配的;
\item[(iii)] 对一切 $\alpha, \beta \in R$, $a \in M$, 有 $(\alpha \cdot \beta) \cdot a = \alpha \cdot (\beta \cdot a)$, 即纯量乘法是结合的;
\item[(iv)] 对单位元 $1 \in R$ 与一切 $a \in M$, 有 $1 \cdot a = a$.
\end{enumerate}
\end{definition}

我们想考虑几个例子:
\begin{enumerate}
\item 域 $K$ 上的向量空间是 $K$-模.
\item 设 $R$ 是带单位元 $1$ 的交换环, $\mathfrak{a} \subset R$ 是理想.则 $R$ 的加法与乘法在 $\mathfrak{a}$ 上诱导出 $R$-模的结构, 特别地 $R$ 自身是 $R$-模.对 $n \in \mathbb{N} - \{0\}$, $R^n$ 以明显的方式是 $R$-模.
\item 每个阿贝尔群 $G$ 都是 $\mathbb{Z}$-模; 如通常那样, 对 $n \in \mathbb{N}$ 和 $g \in G$, 把 $n \cdot g$ 定义为 $g$ 的 $n$ 重和, 并把 $(-n) \cdot g$ 定义为 $-(n \cdot g)$.
\item $K$-向量空间 $V$ 的每个自同态 $f \colon V \longrightarrow V$ 都如上所述在 $V$ 上诱导出 $K[T]$-模的结构.若 $V$ 有限维并且 $p_f$ 表示 $f$ 的最小多项式, 则对一切 $v \in V$ 有 $p_f \cdot v = 0$.与向量空间不同, 对 $R$-模 $M$, 从方程 $\alpha \cdot m = 0$($\alpha \in R$, $m \in M$) 不能推出 $\alpha$ 或 $m$ 消失.
\end{enumerate}

在向量空间中起作用的一大批概念, 对模也有意义.设 $M$ 是带单位元 $1$ 的交换环 $R$ 上的模.$M$ 的子模是子集 $N \subset M$, 非空且满足:

\[
\begin{array}{c}
a, b \in N \Longrightarrow a + b \in N, \\
\alpha \in R, \; a \in N \Longrightarrow \alpha \cdot a \in N
\end{array}\]

则 $N$ 在由 $M$ 继承的运算下又构成 $R$-模.若把 $R$ 看作自身上的模, 则环 $R$ 的理想与 $R$ 的子模一致.

$R$-模 $M$ 与 $N$ 之间的同态是映射 $\varphi \colon M \longrightarrow N$, 满足

\[
\varphi(a + b) = \varphi(a) + \varphi(b), \quad \varphi(\alpha \cdot a) = \alpha \cdot \varphi(a),\]

其中 $a, b \in M$, $\alpha \in R$.此时也说 $R$-线性映射.$R$-模的单态,满态即同构如通常那样定义为单射,满射即双射的同态.

对 $M$ 中元素的系统 $(a_i)_{i \in I}$, 可以考虑由它生成的子模 $M' \subset M$, 后者由

\[
M' := \sum_{i \in I} R a_i := \Big\{\sum_{i \in I} \alpha_i a_i;\; \alpha_i \in R, \;\alpha_i = 0 \text{ 对几乎一切 } i \in I\Big\}\]

定义.若 $M = \sum_{i \in I} R a_i$, 则称 $(a_i)_{i \in I}$ 为 $M$ 的生成系.若 $M$ 有一个有限生成系, 则称 $M$ 有限生成, 或者(用不规范的措辞)有限.此外, $M$ 中元素的系统 $(a_i)_{i \in I}$ 称为自由的(或线性无关的), 如果从方程

\[
\sum_{i \in I} \alpha_i a_i = 0\]

其中系数 $\alpha_i \in R$ 对几乎一切 $i \in I$ 消失, 已经推出对一切 $i$ 有 $\alpha_i = 0$.自由的生成系也称为基.但要注意, 与向量空间不同, 模一般没有基; 参见上面的例 (4).有基的模称为自由的, 若有有限基则称有限自由的.自由 $R$-模之间的同态可以如通常那样关于选取的基用系数在 $R$ 中的矩阵描述.

若 $M_1, \ldots, M_n$ 是 $R$-模 $M$ 的子模, 则可以考虑其和

\[
M' = \Big\{\sum_{i = 1}^{n} a_i;\; a_i \in M_i, \; i = 1, \ldots, n\Big\}\text{.}\]

此时 $M'$ 又是 $M$ 的子模, 记 $M' = \sum_{i = 1}^{n} M_i$.此外, 若对每个 $a \in M'$, 表示 $a = \sum_{i = 1}^{n} a_i$(元素 $a_i \in M_i$) 都是唯一的, 则称 $M'$ 为 $M_i$ 的直和, 记为 $M' = \bigoplus_{i = 1}^{n} M_i$.

与向量空间的情形一样, 未必作为 $R$-模 $N$ 的子模给出的 $R$-模 $N_1, \ldots, N_n$ 的直和 $\bigoplus_{i = 1}^{n} N_i$ 也可以构造.只需令

\[
\bigoplus_{i = 1}^{n} N_i = N_1 \times \ldots \times N_n\]

并把这个笛卡尔积在分量加法即分量纯量乘法下看作 $R$-模.于是对 $i = 1, \ldots, n$, $N_i$ 可以分别与子模

\[
0 \times \ldots \times 0 \times N_i \times 0 \times \ldots \times 0 \subset N_1 \times \ldots \times N_n\]

等同, 使得模 $N_1 \times \ldots \times N_n$ 确实可以看作子模 $N_1, \ldots, N_n$ 的直和.

若 $N$ 是 $R$-模 $M$ 的子模, 则可以作\emph{剩余类模}(\textup{Restklassenmodul})\index{shengyuleimo@剩余类模(Restklassenmodul)} $M/N$.与向量空间的情形一样, $N$ 在 $M$ 上诱导出等价关系:

\[
a \sim b \Longleftrightarrow a - b \in N\]

$M/N$ 由相应的等价类, 即元素 $a \in M$ 的陪集 $[a] = a + N$ 组成.$M/N$ 上的 $R$-模结构由公式

\[
\begin{array}{c c}
[a] + [b] = [a + b], & \quad a, b \in M, \\
\alpha \cdot [a] = [\alpha \cdot a], & \quad \alpha \in R, \; a \in M,
\end{array}\]

给出, 其中当然要检验良定义性.同态定理~\mtref{thm:2.2.8} 与~\mtref{cor:2.2.9} 可以不加修改地搬过来.特别地, $R$-\emph{模同态}(\textup{Homomorphismus von Moduln})\index{motongtai@模同态(Homomorphismus von Moduln)} $\varphi \colon M \longrightarrow N$ 恒诱导出单射的 $R$-模同态 $\bar{\varphi} \colon M/\ker \varphi \longrightarrow N$.

由于模一般没有基, 维数的概念不能随意从向量空间搬到模.维数概念的某些方面由所谓的\emph{模的长度}(\textup{Länge eines Moduls})\index{modeliangdu@模的长度(Länge eines Moduls)}所涵盖.这里对 $R$-模 $M$ 指的是类型为

\[
0 \subsetneq M_1 \subsetneq M_2 \subsetneq \ldots \subsetneq M_\ell = M\]

的子模严格上升链的长度 $\ell$ 的上确界 $\ell_R(M)$.例如 $\ell_R(M) = 0$ 等价于 $M = 0$.此外, $\mathbb{Z}$ 作为自身上的\emph{自由模}(\textup{freier Modul})\index{ziyoumo@自由模(freier Modul)}具有长度 $\infty$.作为后面的工具我们需要两个引理.

\begin{lemma}\label{lem:6.3.2}
设 $R$ 是主理想环, $a = p_1 \ldots p_r$ 是 $R$ 中元素 $a \neq 0$ 的素因子分解.则剩余类模 $R/aR$ 具有长度 $\ell_R(R/aR) = r$.
\end{lemma}

\begin{proof}
设 $\pi \colon R \longrightarrow R/aR$ 是典范投影.由于在对应 $\bar{\mathfrak{a}} \longmapsto \pi^{-1}(\bar{\mathfrak{a}})$ 下, $R/aR$ 的理想 $\bar{\mathfrak{a}} \subset R/aR$ 与 $R$ 中包含 $aR$ 的理想双射对应, $R/aR$ 的长度与类型为

\[
Ra \subsetneq \mathfrak{a}_1 \subsetneq \mathfrak{a}_2 \subsetneq \ldots \subsetneq \mathfrak{a}_\ell = R\]

的严格上升理想链的长度 $\ell$ 的上确界一致.由于 $R$ 是主理想环, 每个 $\mathfrak{a}_i$ 由某个元素 $a_i$ 生成.此外, 真包含 $\mathfrak{a}_{i - 1} \subsetneq \mathfrak{a}_i$ 等价于 $a_i$ 是 $a_{i - 1}$ 的真因子.因此 $R/aR$ 的长度等于一切 $\ell \in \mathbb{N}$ 的上确界, 其中要求存在 $a_1, \ldots, a_\ell \in R$ 使 $a_i$ 分别是 $a_{i - 1}$ 的真因子, $i = 1, \ldots, \ell$; 这里应令 $a_0 = a$.由于在 $R$ 中成立唯一素因子分解定理, 并且 $a_0$ 是 $r$ 个素因子的乘积, 这个上确界算出为 $r$.
\end{proof}

\begin{lemma}\label{lem:6.3.3}
若 $R$-模 $M$ 是两个子模 $M'$ 与 $M''$ 的直和, 则 $\ell_R(M) = \ell_R(M') + \ell_R(M'')$.
\end{lemma}

\begin{proof}
若有子模的严格上升链

\[
0 \subsetneq M_1' \subsetneq M_2' \subsetneq \ldots \subsetneq M_r' = M',\]

\[
0 \subsetneq M_1'' \subsetneq M_2'' \subsetneq \ldots \subsetneq M_s'' = M'',\]

则

\[
\begin{array}{c}
0 \subsetneq M_1' \oplus 0 \subsetneq M_2' \oplus 0 \subsetneq \ldots \subsetneq M_r' \oplus 0 \\
\subsetneq M_r' \oplus M_1'' \subsetneq M_r' \oplus M_2'' \subsetneq \ldots \subsetneq M_r' \oplus M_s'' = M
\end{array}\]

是 $M$ 中长度为 $r + s$ 的严格上升链.于是

\[
\ell_R(M) \geq \ell_R(M') + \ell_R(M'')\text{.}\]

为证明反向估计, 考虑子模的严格上升链

\[
0 = M_0 \subsetneq M_1 \subsetneq M_2 \subsetneq \ldots \subsetneq M_\ell = M\text{.}\]

设 $\pi'' \colon M' \oplus M'' \longrightarrow M''$ 是到第二个直和项的投影, 于是 $\ker \pi'' = M'$.则对 $0 \leq \lambda < \ell$, 如我们马上会看到的, 分别有 $M_\lambda \cap M'$ 真含于 $M_{\lambda + 1} \cap M'$, 或 $\pi''(M_\lambda)$ 真含于 $\pi''(M_{\lambda + 1})$.由此推出 $\ell \leq \ell_R(M') + \ell_R(M'')$, 从而如所希望

\[
\ell_R(M) = \ell_R(M') + \ell_R(M'')\text{.}\]

为证明刚才断言的包含关系, 我们假设 $M_\lambda \cap M' = M_{\lambda + 1} \cap M'$ 且 $\pi''(M_\lambda) = \pi''(M_{\lambda + 1})$, 并证明这已经蕴含 $M_\lambda = M_{\lambda + 1}$, 与我们的假设矛盾.事实上, 对 $a \in M_{\lambda + 1}$, 由于 $\pi''(M_\lambda) = \pi''(M_{\lambda + 1})$, 存在 $a' \in M_\lambda$ 使 $\pi(a) = \pi(a')$, 即 $a - a' \in \ker \pi = M'$.则由 $a, a' \in M_{\lambda + 1}$ 甚至得到 $a - a' \in M_{\lambda + 1} \cap M' = M_\lambda \cap M'$, 从而 $a = a' + (a - a') \in M_\lambda$.因此 $M_{\lambda + 1} \subset M_\lambda$, 即 $M_{\lambda + 1} = M_\lambda$, 但这已被排除.
\end{proof}


\begin{bigtheorem}[\emph{初等因子定理}\index{chudengyinzidingli@初等因子定理(Elementarteilersatz)}]\label{bigthm:6.3.4}\index{chudengyinzidingli@初等因子定理(Elementarteilersatz)}
设 $R$ 是主理想环, $F$ 是有限自由 $R$-模.再设 $M \subset F$ 是子模.则存在 $F$ 的基的一部分 $x_1, \ldots, x_s \in F$ 以及系数 $\alpha_1, \ldots, \alpha_s \in R - \{0\}$, 使得:
\begin{enumerate}
\item[(i)] $\alpha_1 x_1, \ldots, \alpha_s x_s$ 构成 $M$ 的一个基;
\item[(ii)] 对 $1 \leq i < s$ 有 $\alpha_i \mid \alpha_{i + 1}$.
\end{enumerate}
其中 $\alpha_1, \ldots, \alpha_s$ 到差一个相伴(即到差一个单位)为止由 $M$ 唯一确定, 与 $x_1, \ldots, x_s$ 的选取无关.称 $\alpha_1, \ldots, \alpha_s$ 为 $M \subset F$ 的初等因子.特别地, 其个数 $s$ 唯一确定.
\end{bigtheorem}

\begin{proof}[定理~\mtref{bigthm:6.3.4} 存在性断言的证明]
设 $Y = (y_1, \ldots, y_m)$ 是 $F$ 的一个基.首先对 $m$ 归纳证明子模 $M \subset F$ 有限生成.$m = 1$ 时这是清楚的, 因为 $F$ 作为 $R$-\emph{模同构}(\textup{Isomorphismus von Moduln})\index{motonggou@模同构(Isomorphismus von Moduln)}于 $R$, 而 $M$ 对应于 $R$ 中的一个理想.后者有限生成, 因为 $R$ 是主理想环.设 $m > 1$.令 $F' = \sum_{i = 1}^{m - 1} R y_i$ 和 $F'' = R y_m$.再考虑投影 $\pi \colon F \longrightarrow F''$, 它把 $i < m$ 的 $y_i$ 映到 $0$, 把 $y_m$ 映到 $y_m$; 则 $\ker \pi = F'$, 并且用 \ref{sec:2.3} 节的术语说, 有短正合序列

\[
0 \longrightarrow F' \longrightarrow F \longrightarrow F'' \longrightarrow 0\text{.}\]

按归纳假设, 子模 $M \cap F' \subset F'$ 与 $\pi(M) \subset F''$ 有限生成, 并且如通常那样(例如如同定理~\mtref{thm:2.1.10} 证明中) 证明 $M \cap F'$ 的一个生成系连同 $\pi(M)$ 的一个生成系的提升共同构成 $M$ 的生成系.故 $M$ 有限生成.

我们保留 $Y = (y_1, \ldots, y_m)$ 作为 $F$ 的基, 并选取 $M$ 的一个有限生成系 $z_1, \ldots, z_n$.若 $e = (e_1, \ldots, e_n)$ 是 $R$-模 $R^n$ 的典范基, 则可以考虑由 $e_j \longmapsto z_j$ 定义的 $R$-线性映射 $f \colon R^n \longrightarrow F$, 其像给出 $M$.若

\[
z_j = \sum_{i = 1}^{m} \alpha_{ij} y_i, \qquad j = 1, \ldots, n,\]

则 $A = (\alpha_{ij})_{i, j} \in R^{m \times n}$ 是 $f$ 关于基 $e$ 和 $Y$ 的矩阵.我们现在使用下面这个稍后将证明的辅助结果:

\begin{lemma}\label{lem:6.3.5}
设 $R$ 是主理想环, $A = (\alpha_{ij}) \in R^{m \times n}$ 是系数在 $R$ 中的矩阵.则存在可逆矩阵 $S \in R^{m \times m}$ 和 $T \in R^{n \times n}$ 使

\[
S \cdot A \cdot T = \left(\begin{array}{cccccccc}
\alpha_1 & 0 & & & 0 & 0 & \ldots & 0 \\
0 & \alpha_2 & & & 0 & 0 & \ldots & 0 \\
 & & \ddots & & & & & \\
 & & & \ddots & & & & \\
0 & 0 & & & \alpha_s & 0 & \ldots & 0 \\
0 & 0 & \ldots & \ldots & 0 & 0 & \ldots & 0 \\
\vdots & \vdots & & & \vdots & \vdots & & \vdots \\
0 & 0 & \ldots & \ldots & 0 & 0 & \ldots & 0
\end{array}\right)\]

且系数 $\alpha_1, \ldots, \alpha_s \in R - \{0\}$(其中 $0 \leq s \leq \min(m, n)$) 满足对 $1 \leq i < s$ 的整除条件 $\alpha_i \mid \alpha_{i + 1}$.其中 $\alpha_1, \ldots, \alpha_s$ 到差一个相伴为止唯一确定; 称它们为矩阵 $A$ 的初等因子.
\end{lemma}

把 $S$ 和 $T$ 理解为基变换矩阵, 则矩阵 $SAT$ 也描述映射 $f$, 不过是关于另外的合适基, 即 $R^n$ 的基 $e_1', \ldots, e_n'$ 和 $F$ 的基 $x_1, \ldots, x_m$.特别地推出, $M$ 作为 $f$ 的像由 $\alpha_1 x_1, \ldots, \alpha_s x_s$ 生成.由于系统 $x_1, \ldots, x_m$ 是自由的, 并且我们考虑整环 $R$ 中的系数, $\alpha_1 x_1, \ldots, \alpha_s x_s$ 甚至构成 $M$ 的一个基.这样我们就得到了初等因子 $\alpha_1, \ldots, \alpha_s$ 的 $M \subset F$ 的存在性, 化归为引理~\mtref{lem:6.3.5} 的存在性断言.
\end{proof}

\begin{proof}[引理~\mtref{lem:6.3.5} 存在性断言的证明]
我们首先设 $R$ 是欧几里得环, 并用带余除法给出一个构造性方法, 证明矩阵 $A = (\alpha_{ij})$ 可以通过可逆的初等行变换与列变换化成所要求的形状, 即通过交换行(或列)以及把某行(或某列)的倍数加到另一行(或另一列).与系数在一个域中的矩阵情形一样, 这种类型的初等变换可解释为左乘(或右乘)一个可逆的初等矩阵.所需的行变换因此总体上对应于左乘一个可逆矩阵 $S \in R^{m \times m}$, 所需的列变换相应地对应对右乘一个可逆矩阵 $T \in R^{n \times n}$.然后我们把方法推广, 使经过修改也能用于主理想环.

因此下面首先考虑带次数映射 $\delta \colon R - \{0\} \longrightarrow \mathbb{N}$ 的欧几里得环 $R$.对 $A = 0$ 无需证明, 所以可设 $A \neq 0$.我们的策略是通过初等变换逐步减小最小值

\[
d(A) := \min\{\delta(\alpha);\; \alpha \text{ 是 } A \text{ 的非零系数}\}\text{.}\]

由于 $\delta$ 在 $\mathbb{N}$ 中取值, 这个过程经过有限步后必须中止.若此时 $\alpha \neq 0$ 是变换后的矩阵中次数 $\delta(\alpha)$ 最小的系数, 则我们用带余除法证明 $\alpha$ 整除矩阵的一切其他系数; $\alpha$ 就是 $A$ 的第一个初等因子.

具体步骤如下.通过交换 $A$ 的行与列, 可以假设 $d(A) = \delta(\alpha_{11})$, 即 $\delta(\alpha_{11})$ 在一切 $\delta(\alpha_{ij})$(其中 $\alpha_{ij} \neq 0$) 中最小.在每一步的开头我们都建立这种情形.若第 1 列中的某个元素, 比如 $\alpha_{i1}$, 不被 $\alpha_{11}$ 整除, 则用 $\alpha_{11}$ 去除 $\alpha_{i1}$, 比如 $\alpha_{i1} = q \alpha_{11} + \beta$, 其中 $\delta(\beta) < \delta(\alpha_{11})$, 并把第 1 行的 $q$ 倍从第 $i$ 行减去.结果在位置 $(i, 1)$ 出现元素 $\beta$.$A$ 的非零系数的次数的最小值 $d(A)$ 因此已经减小, 于是用下一步重新开始这个过程.同样, 可以用初等列变换修改第 1 行的元素.由于 $d(A)$ 在 $\mathbb{N}$ 中取值, 不能无限次地减小, 经过有限步后第 1 列以及第 1 行的每个元素都是 $\alpha_{11}$ 的倍数, 并且可以通过把第 1 行的倍数加到矩阵的其余行, 假设对 $i > 1$ 有 $\alpha_{i1} = 0$.相应地可以处理第 1 行, 并这样达到对 $i, j > 1$ 有 $\alpha_{i1} = \alpha_{1j} = 0$.这里我们还可以假设最小值 $d(A)$ 与 $\delta(\alpha_{11})$ 一致; 否则要重新开始过程, 并把相应的元素重新放到位置 $(1, 1)$.若现在存在 $i, j > 1$ 使 $\alpha_{11} \nmid \alpha_{ij}$, 则把第 $j$ 列加到第 1 列, 这个过程保持 $\alpha_{11}$ 不变.如刚才所述, $\alpha_{11}$ 下方的元素可以重新化零, 同时减小次数 $d(A)$.经过有限步这样得到矩阵 $(\alpha_{ij})$, 满足对 $i, j > 1$ 有 $\alpha_{i1} = \alpha_{1j} = 0$, 并且 $\alpha_{11}$ 整除一切 $i, j > 1$ 的其他元素 $\alpha_{ij}$.然后以同样的方式处理 $A = (\alpha_{ij})$ 的子矩阵 $(\alpha_{ij})_{i, j > 1}$, 只要它还不是零矩阵.为此所需的变换保持 $A$ 的第一行与第一列不变, 特别保持 $\alpha_{11}$ 整除 $A$ 的其余一切系数这一条件.把这个过程归纳地继续下去, 经过有限步最终得到一个矩阵, 在其主对角线上依次排列着所求的初等因子并具有断言的整除性, 而其余一切元素都消失.这样就证明了引理~\mtref{lem:6.3.5} 的存在性断言, 特别也证明了定理~\mtref{bigthm:6.3.4} 的存在性断言, 至少对欧几里得环 $R$ 的情形.
若只知道 $R$ 是主理想环, 则需要稍作推广的初等矩阵变换, 我们先说明这一点.设 $\sigma, \tau, \sigma', \tau' \in R$ 满足 $\sigma \tau' - \tau \sigma' = 1$.则矩阵

\[
\left(\begin{array}{cc} \sigma & \tau \\ \sigma' & \tau' \end{array}\right), \qquad \left(\begin{array}{cc} \tau' & -\tau \\ -\sigma' & \sigma \end{array}\right) \in R^{2 \times 2}\]

互为逆, 特别可逆.相应地看到, 对 $\sigma \tau' - \tau \sigma' = \pm 1$ 和 $1 \leq i < j \leq m$, 类型为

\[
E^{ij}(\sigma, \tau, \sigma', \tau') = \left(\begin{array}{ccccc}
1 & & & & \\
 & \sigma & & \tau & \\
 & & 1 & & \\
 & \sigma' & & \tau' & \\
 & & & & 1
\end{array}\right) \in R^{m \times m}\]

的矩阵也可逆.这里元素 $\sigma, \tau, \sigma', \tau'$ 分别位于位置 $(i, i)$, $(i, j)$, $(j, i)$, $(j, j)$, 而对角线上的"$1$"表示一串元素 $1$.此外, 矩阵 $E^{ij}(\sigma, \tau, \sigma', \tau')$ 的其他位置用 $0$ 填充, 出于清晰起见没有写出.若把 $A$ 从左乘 $E^{ij}(\sigma, \tau, \sigma', \tau')$, 则有如下效果: 新的第 $i$ 行是旧的第 $i$ 行的 $\sigma$ 倍与旧的第 $j$ 行的 $\tau$ 倍之和; 相应地, 新的第 $j$ 行是旧的第 $i$ 行的 $\sigma'$ 倍与旧的第 $j$ 行的 $\tau'$ 倍之和.例如, 用常数

\[
\sigma = 0, \qquad \tau = 1, \qquad \sigma' = 1, \qquad \tau' = 0\]

得到第 $i$ 行与第 $j$ 行的交换, 而用

\[
\sigma = 1, \qquad \tau = \varepsilon, \qquad \sigma' = 0, \qquad \tau' = 1\]

得到把第 $j$ 行的 $\varepsilon$ 倍加到第 $i$ 行.通过转置看到, $A$ 的相应列变换可由右乘类型为 $E^{ij}(\sigma, \tau, \sigma', \tau') \in R^{n \times n}$ 的矩阵生成.

现在对元素 $\alpha \in R - \{0\}$, 用 $\delta(\alpha)$ 表示 $\alpha$ 的素因子的个数; 按引理~\mtref{lem:6.3.2}, 这正好是剩余类环 $R/\alpha R$ 的长度.此外, 类似于欧几里得环的情形, 令

\[
d(A) := \min\{\delta(\alpha);\; \alpha \text{ 是 } A \text{ 的非零系数}\}\text{.}\]

其目标是逐步减小 $d(A)$, 直到 $A$ 存在一个整除其余一切系数的系数 $\alpha$.通过交换行与列可以再次假设 $d(A) = \delta(\alpha_{11})$.若第 1 列的某个元素, 比如 $\alpha_{i1}$, 不是 $\alpha_{11}$ 的倍数, 则作 $\alpha_{11}$ 与 $\alpha_{i1}$ 的最大公因子 $\beta$.对这个 $\beta$ 必有 $\delta(\beta) < \delta(\alpha_{11})$, 并且按定理~\mtref{thm:5.2.16}, $\beta$ 生成理想 $R\alpha_{11} + R\alpha_{i1}$, 即存在类型为

\[
\beta = \sigma \alpha_{11} + \tau \alpha_{i1}\]

的方程, 其中 $\sigma, \tau \in R$ 必然互素, 从而满足某个类型为

\[
\sigma \tau' - \tau \sigma' = 1\]

的方程, 其中 $\sigma', \tau' \in R$ 是某些元素.若现在把 $A$ 从左乘 $E^{1i}(\sigma, \tau, \sigma', \tau')$, 则这个过程在 $A$ 的位置 $(1, 1)$ 放置元素 $\beta$, 从而减小最小值 $d(A)$.像欧几里得环的情形那样迭代这个过程, 最终可以达到: 元素 $\alpha_{21}, \ldots, \alpha_{m1}$ 都能被 $\alpha_{11}$ 整除, 或者通过从其余行减去第 1 行的适当倍数, 使得 $\alpha_{21} = \ldots = \alpha_{m1} = 0$ 成立.同样可以用相应的列变换化零元素 $\alpha_{12}, \ldots, \alpha_{1n}$, 等等.因此我们看到, 矩阵 $A$ 可以像欧几里得环的情形那样逐步变换, 直到最终达到所要求的形状.
\end{proof}

我们想用一个简单的例子演示为欧几里得环描述的方法, 为此考虑矩阵

\[
A = \left(\begin{array}{ccc} 6 & 2 & 5 \\ 32 & 2 & 28 \\ 30 & 2 & 26 \end{array}\right) \in \mathbb{Z}^{3 \times 3}\text{.}\]

为方便起见, 在确定 $A$ 的初等因子时, 除上面使用的初等行变换与列变换外, 还允许把某行或某列乘以环 $R = \mathbb{Z}$ 的一个单位.这是允许的, 因为这种变换也可以解释为左乘或右乘可逆的初等矩阵, 即把单位矩阵的某个对角元素 $1$ 换成 $R$ 中的一个单位所得的那种矩阵.否则我们还是坚持为欧几里得环描述的方法, 尽管通过更巧妙选择初等变换可以简化 $A$ 的初等因子的确定.

\[
A = \left(\begin{array}{ccc} 6 & 2 & 5 \\ 32 & 2 & 28 \\ 30 & 2 & 26 \end{array}\right)
\longmapsto^{(1)} \left(\begin{array}{ccc} 2 & 6 & 5 \\ 2 & 32 & 28 \\ 2 & 30 & 26 \end{array}\right)
\longmapsto^{(2)} \left(\begin{array}{ccc} 2 & 6 & 5 \\ 0 & 26 & 23 \\ 0 & 24 & 21 \end{array}\right)\]

\[
\longmapsto^{(3)} \left(\begin{array}{ccc} 2 & 0 & 1 \\ 0 & 26 & 23 \\ 0 & 24 & 21 \end{array}\right)
\longmapsto^{(4)} \left(\begin{array}{ccc} 1 & 0 & 2 \\ 23 & 26 & 0 \\ 21 & 24 & 0 \end{array}\right)
\longmapsto^{(5)} \left(\begin{array}{ccc} 1 & 0 & 2 \\ 0 & 26 & -46 \\ 0 & 24 & -42 \end{array}\right)\]

\[
\longmapsto^{(6)} \left(\begin{array}{ccc} 1 & 0 & 0 \\ 0 & 26 & 46 \\ 0 & 24 & 42 \end{array}\right)
\longmapsto^{(7)} \left(\begin{array}{ccc} 1 & 0 & 0 \\ 0 & 24 & 42 \\ 0 & 26 & 46 \end{array}\right)
\longmapsto^{(8)} \left(\begin{array}{ccc} 1 & 0 & 0 \\ 0 & 24 & 42 \\ 0 & 2 & 4 \end{array}\right)\]

\[
\longmapsto^{(9)} \left(\begin{array}{ccc} 1 & 0 & 0 \\ 0 & 2 & 4 \\ 0 & 24 & 42 \end{array}\right)
\longmapsto^{(10)} \left(\begin{array}{ccc} 1 & 0 & 0 \\ 0 & 2 & 4 \\ 0 & 0 & -6 \end{array}\right)
\longmapsto^{(11)} \quad \left(\begin{array}{ccc} 1 & 0 & 0 \\ 0 & 2 & 0 \\ 0 & 0 & 6 \end{array}\right)\]

于是得到 $1, 2, 6$ 为 $A$ 的初等因子, 其中具体执行了下列初等变换:
\begin{enumerate}
\item 交换第 1 列与第 2 列;
\item 从第 2 行与第 3 行减去第 1 行;
\item 从第 2 列与第 3 列分别减去第 1 列的 $3$ 倍与 $2$ 倍;
\item 交换第 1 列与第 3 列;
\item 从第 2 行与第 3 行分别减去第 1 行的 $23$ 倍与 $21$ 倍;
\item 从第 3 列减去第 1 列的 $2$ 倍, 并把第 3 列乘以 $-1$;
\item 交换第 2 行与第 3 行;
\item 从第 3 行减去第 2 行;
\item 交换第 2 行与第 3 行;
\item 从第 3 行减去第 2 行的 $12$ 倍;
\item 从第 3 列减去第 2 列的 $2$ 倍, 并把第 3 列乘以 $-1$.
\end{enumerate}

接下来我们处理引理~\mtref{lem:6.3.5} 即定理~\mtref{bigthm:6.3.4} 中的唯一性断言, 从一个基本引理开始, 它在后面也有意义.

\begin{lemma}\label{lem:6.3.6}
设 $R$ 是主理想环, 并且

\[
Q \simeq \bigoplus_{i = 1}^{s} R/\alpha_i R\]

是 $R$-模的同构, 其中 $\alpha_1, \ldots, \alpha_s \in R - \{0\}$ 是非单位, 对 $1 \leq i < s$ 有 $\alpha_i \mid \alpha_{i + 1}$, 而 $\bigoplus_{i = 1}^{s} R/\alpha_i R$ 表示 $R$-模 $R/\alpha_i R$ 的构造性直和.则 $\alpha_1, \ldots, \alpha_s$ 到差一个相伴为止由 $Q$ 唯一确定.
\end{lemma}

\begin{proof}
出于技术原因, 我们反转 $\alpha_i$ 的编号, 并考虑两个分解

\[
Q \simeq \bigoplus_{i = 1}^{s} R/\alpha_i R \simeq \bigoplus_{j = 1}^{t} R/\beta_j R\]

其中对 $1 \leq i < s$ 有 $\alpha_{i + 1} \mid \alpha_i$, 对 $1 \leq j < t$ 有 $\beta_{j + 1} \mid \beta_j$.若存在指标 $k \leq \min\{s, t\}$ 使 $\alpha_k R \neq \beta_k R$, 则选取具有此性质的最小 $k$.由于对 $1 \leq i < k$ 有 $\alpha_i R = \beta_i R$, 并且 $\alpha_{k + 1}, \ldots, \alpha_s$ 全部是 $\alpha_k$ 的因子, $\alpha_k Q$ 分解为

\[
\begin{array}{c}
\alpha_k Q \simeq \bigoplus_{i = 1}^{k - 1} \alpha_k \cdot (R/\alpha_i R) \\
\simeq \bigoplus_{i = 1}^{k - 1} \alpha_k \cdot (R/\alpha_i R) \oplus \bigoplus_{j = k}^{t} \alpha_k \cdot (R/\beta_j R)\text{.}
\end{array}\]

我们利用引理~\mtref{lem:6.3.2} 和引理~\mtref{lem:6.3.3}.由于 $\ell_R(\alpha_k Q) \leq \ell_R(Q) < \infty$, 比较两边得到对 $j = k, \ldots, t$ 有 $\ell_R(\alpha_k \cdot (R/\beta_j R)) = 0$.后者特别意味着 $\alpha_k \cdot (R/\beta_k R) = 0$, 即 $\alpha_k R \subset \beta_k R$.相应地证明 $\beta_k R \subset \alpha_k R$, 从而 $\alpha_k R = \beta_k R$, 与我们的假设矛盾.因此对一切满足 $1 \leq i \leq \min\{s, t\}$ 的指标 $i$ 有 $\alpha_i R = \beta_i R$.若再有 $s \leq t$, 则再次利用引理~\mtref{lem:6.3.3} 推出 $\bigoplus_{j = s + 1}^{t} R/\beta_j R$ 的长度为 $0$, 即对 $j = s + 1, \ldots, t$, $R/\beta_j R$ 的长度为 $0$.另一方面, 对 $j = 1, \ldots, t$, $\beta_j$ 不是单位, 因此按引理~\mtref{lem:6.3.2}, 每个模 $R/\beta_j R$ 都有大于 $0$ 的长度, 于是得到 $s = t$.
\end{proof}

\begin{proof}[定理~\mtref{bigthm:6.3.4} 与引理~\mtref{lem:6.3.5} 唯一性断言的证明]
引理~\mtref{lem:6.3.5} 中的唯一性是定理~\mtref{bigthm:6.3.4} 中唯一性断言的推论, 因此只需证明定理~\mtref{bigthm:6.3.4} 中初等因子的唯一性.设 $x_1, \ldots, x_s$ 是 $F$ 的基的一部分, $\alpha_1, \ldots, \alpha_s \in R$ 满足对 $1 \leq i < s$ 有 $\alpha_i \mid \alpha_{i + 1}$, 使得 $\alpha_1 x_1, \ldots, \alpha_s x_s$ 构成 $M$ 的基.则子模

\[
F' = \bigoplus_{i = 1}^{s} R x_i = \{a \in F;\; \text{存在 } \alpha \in R - \{0\} \text{ 使 } \alpha a \in M\}\]

只依赖于 $M$, 而不依赖于元素 $x_1, \ldots, x_s$ 的特殊选取.典范同态

\[
F' \longrightarrow \bigoplus_{i = 1}^{s} R/\alpha_i R, \quad \sum_{i = 1}^{s} \gamma_i x_i \longmapsto (\bar{\gamma}_1, \ldots, \bar{\gamma}_s),\]

其中 $\bar{\gamma}_i$ 分别是 $\gamma_i$ 在 $R/R\alpha_i$ 中的剩余类, 是满射并以 $M$ 为核, 因此按同态定理诱导同构

\[
F'/M \xrightarrow{\sim} \bigoplus_{i = 1}^{s} R/\alpha_i R\text{.}\]

由引理~\mtref{lem:6.3.6} 的唯一性断言可以推出, $\alpha_1, \ldots, \alpha_s$ 中的非单位至少到差一个相伴为止由 $M$ 唯一确定.

为了看出 $\alpha_1, \ldots, \alpha_s$ 中单位的个数也由 $M$ 唯一确定, 我们想证明 $s$ 作为 $M$ 即 $F'$ 的基的元素个数唯一确定.由于 $M$ 的每个基 $z_1, \ldots, z_s$ 都给出同构 $M \longrightarrow R^s$, 只需证明同构 $R^s \longrightarrow R^{s'}$ 的存在已经蕴含 $s = s'$.在 $R$ 是域的情形, 这由向量空间的维数理论得出; 参见例如定理~\mtref{thm:2.1.8}.但若 $R$ 不是域, 则 $R - R^*$ 至少含有一个非零元素, 而按推论~\mtref{cor:5.2.13} 这个元素可以写成素元素的乘积.因此在 $R$ 中至少可以找到一个素元素 $p$, 并且容易验证, $R$-模的每个同构 $R^s \longrightarrow R^{s'}$ 都诱导出 $R/pR$-模的同构

\[
(R/pR)^s = R^s/pR^s \xrightarrow{\sim} R^{s'}/pR^{s'} = (R/pR)^{s'}\text{.}\]

但按推论~\mtref{cor:5.2.17}, $R/pR$ 是域, 于是又可以推出 $s = s'$.此外也可以利用引理~\mtref{lem:6.3.6} 的唯一性断言.
\end{proof}

\begin{corollary}\label{cor:6.3.7}
设 $R$ 是主理想环, $F$ 是有限自由 $R$-模.则 $F$ 的任意两个基具有相同的长度.这个长度也称为 $F$ 的秩.
\end{corollary}

\begin{proof}
这个断言已经在上述证明中推导出来, 用来验证子模 $M \subset F$ 的初等因子的个数唯一确定.此外, 推论~\mtref{cor:6.3.7} 的断言也形式地由定理~\mtref{bigthm:6.3.4} 的唯一性断言得出, 因为平凡子模 $F \subset F$ 的初等因子的个数唯一确定.
\end{proof}

对于定理~\mtref{bigthm:6.3.4} 情形中初等因子的计算, 重要的是知道我们把这些初等因子作为某个矩阵的初等因子得到, 因为矩阵的初等因子可以用引理~\mtref{lem:6.3.5} 证明中给出的实用方法确定.我们想把精确的条件在这里再单独表述一下.

\begin{corollary}\label{cor:6.3.8}
设 $R$ 是主理想环, $F$ 是有限自由 $R$-模, $M \subset F$ 是子模.则 $M$ 有限生成.

设 $F'$ 是另一个有限自由 $R$-模, $f \colon F' \longrightarrow F$ 是 $R$-线性映射, 把 $F'$ 的一个基映到 $M$ 的一个生成系, 即具有像 $\operatorname{im} f = M$.若 $A$ 是系数在 $R$ 中,关于 $F'$ 与 $F$ 中合适基描述 $f$ 的矩阵, 则 $A$ 的初等因子与 $M \subset F$ 的初等因子一致.
\end{corollary}

\aufgaben

除另有要求外, 设 $R$ 是带单位元 $1$ 的交换环.

\noindent\textbf{习题 1.} 确定下列矩阵的初等因子:

\[
\left(\begin{array}{ccc} 2 & 6 & 8 \\ 3 & 1 & 2 \\ 9 & 5 & 4 \end{array}\right), \qquad \left(\begin{array}{ccc} 4 & 0 & 0 \\ 0 & 10 & 0 \\ 0 & 0 & 15 \end{array}\right) \in \mathbb{Z}^{3 \times 3}\]

\medskip
\noindent\textbf{习题 2.} 设 $R$ 是主理想环, $A = (\alpha_{ij})_{i, j} \in R^{m \times n}$ 是系数在 $R$ 中的非平凡矩阵.若 $\alpha_1, \ldots, \alpha_s$ 是 $A$ 的初等因子, 则 $s > 0$ 且

\[
\alpha_1 = \mathrm{ggT}(\alpha_{ij};\; i = 1, \ldots, m, \; j = 1, \ldots, n)\text{.}\]

\medskip
\noindent\textbf{习题 3.} 设 $a_{11}, \ldots, a_{1n}$ 是主理想环 $R$ 中互素的元素, 即 $\mathrm{ggT}(a_{11}, \ldots, a_{1n}) = 1$.证明: 存在元素 $a_{ij} \in R$, $i = 2, \ldots, n$, $j = 1, \ldots, n$, 使矩阵 $(a_{ij})_{i, j = 1, \ldots, n}$ 在 $R^{n \times n}$ 中可逆.(AT 352)

\medskip
\noindent\textbf{习题 4.} 设 $f \colon M \longrightarrow N$ 是主理想环 $R$ 上有限自由模(即 $M$ 与 $N$ 各有有限基) 之间的同态.利用初等因子定理的断言, 推出存在 $M$ 的基 $X$ 与 $N$ 的基 $Y$, 以及非零元素 $\alpha_1, \ldots, \alpha_s \in R$ 满足对 $1 \leq i < s$ 有 $\alpha_i \mid \alpha_{i + 1}$, 使得

\[
A_{f, X, Y} = \left(\begin{array}{cccccccc}
\alpha_1 & 0 & & & 0 & 0 & \ldots & 0 \\
0 & \alpha_2 & & & 0 & 0 & \ldots & 0 \\
 & & \ddots & & & & & \\
 & & & \ddots & & & & \\
0 & 0 & \ldots & \ldots & \alpha_s & 0 & \ldots & 0 \\
0 & 0 & \ldots & \ldots & 0 & 0 & \ldots & 0 \\
\vdots & \vdots & & & \vdots & \vdots & & \vdots \\
0 & 0 & \ldots & \ldots & 0 & 0 & \ldots & 0
\end{array}\right)\text{.}\]

其中元素 $\alpha_1, \ldots, \alpha_s$ 到差一个相伴为止唯一确定.

\medskip
\noindent\textbf{习题 5.} 确定 $(\mathbb{Z}/15\mathbb{Z})^4$ 作为 $\mathbb{Z}$-模的长度.

\medskip
\noindent\textbf{习题 6.} 设 $M$ 是有限生成且自由的 $R$-模.证明(AT 353):
\begin{enumerate}
\item[(i)] $M$ 有一个有限基;
\item[(ii)] $M$ 的任意两个基由同样多的元素组成.
\end{enumerate}

\medskip
\noindent\textbf{习题 7.} 在初等因子定理(定理~\mtref{bigthm:6.3.4}) 中我们考虑的是有限自由模 $F$ 的子模 $M$.证明: 若改用 $F$ 自由而 $M$ 有限生成的假设, 定理的断言仍然成立.

\medskip
\noindent\textbf{习题 8.} 设 $V$ 是域 $K$ 上的向量空间.对每个自同态 $f \colon V \longrightarrow V$, 可以在 $V$ 上考虑相应的 $K[T]$-模结构, 它由对元素 $v \in V$ 的 $T \cdot v = f(v)$ 刻画.证明: 用这种方式在 $\operatorname{End}_K(V)$ 与 $V$ 上与 $V$ 作为 $K$-向量空间的结构相容的 $K[T]$-模结构集之间得到双射.

\medskip
\noindent\textbf{习题 9.}(模的第一同构定理) 设 $N, N'$ 是 $R$-模 $M$ 的子模.证明: 典范映射 $N \longrightarrow N + N' \longrightarrow (N + N')/N'$ 以 $N \cap N'$ 为核, 并诱导同构

\[
N/(N \cap N') \xrightarrow{\sim} (N + N')/N'\text{.}\]

\medskip
\noindent\textbf{习题 10.}(模的第二同构定理) 设 $M$ 是 $R$-模, $N \subset N' \subset M$ 是子模.证明:
\begin{enumerate}
\item[(i)] 典范映射 $N' \longrightarrow M \longrightarrow M/N$ 以 $N$ 为核, 并诱导单态 $N'/N \longrightarrow M/N$.因此可以把 $N'/N$ 与其在 $M/N$ 中的像等同, 从而看作 $M/N$ 的子模;
\item[(ii)] 投影 $M \longrightarrow M/N'$ 穿过 $M/N$, 即可以写成复合 $M \xrightarrow{\pi} M/N \xrightarrow{f} M/N'$, 其中 $f$ 是某个模同态, $\pi$ 是典范投影;
\item[(iii)] $f$ 以 $N'/N$ 为核, 并诱导同构

\[
(M/N)/(N'/N) \xrightarrow{\sim} M/N'\text{.}\]
\end{enumerate}


\section{主理想环上有限生成的模}\index{mo@模(Modul)}\index{zhulixianghuan@主理想环(Hauptidealring)}\label{sec:6.4}

我们现在想从初等因子定理推出主理想环上\emph{有限生成模}(\textup{endlich erzeugter Modul})\index{youxianshengchengmo@有限生成模(endlich erzeugter Modul)}的结构的一些推论.在下一节中, 要把得到的结构定理转化为关于向量空间的自同态的标准形的结果.作为工具还需要所谓\emph{中国剩余定理}(\textup{Chinesischer Restsatz})\index{zhongguoshengyudingli@中国剩余定理(Chinesischer Restsatz)}, 我们首先证明它.

\begin{theorem}[中国剩余定理]\label{thm:6.4.1}\index{zhongguoshengyudingli@中国剩余定理(Chinesischer Restsatz)}
设 $R$ 是主理想环, 并且

\[
a = \varepsilon p_1^{n_1} \ldots p_r^{n_r}\]

是 $R$ 中的素因子分解, 其中 $\varepsilon$ 是单位, $p_i$ 是两两不相伴的素元素.若

\[
\pi_i \colon R \longrightarrow R/(p_i^{n_i}), \qquad i = 1, \ldots, r,\]

分别是典范投影, 则同态

\[
\varphi \colon R \longrightarrow R/(p_1^{n_1}) \times \ldots \times R/(p_r^{n_r}), \quad b \longmapsto (\pi_1(b), \ldots, \pi_r(b)),\]

是满射, 并且 $\ker \varphi = (a)$, 从而诱导同构

\[
R/(a) \xrightarrow{\sim} R/(p_1^{n_1}) \times \ldots \times R/(p_r^{n_r})\text{.}\]

其中 $R/(p_1^{n_1}) \times \ldots \times R/(p_r^{n_r})$ 在分量加法与分量乘法下理解为环.
\end{theorem}

\begin{proof}
首先证明 $\varphi$ 是满射.为此显然只需验证存在元素 $e_1, \ldots, e_r \in R$ 使

\[
\pi_i(e_j) = \left\{\begin{array}{l l} 1 & \text{对 } i = j, \\ 0 & \text{否则\text{.}} \end{array}\right.\]

由于对固定的 $j \in \{1, \ldots, r\}$, 元素 $p_j^{n_j}$ 与 $\prod_{i \neq j} p_i^{n_i}$ 互素, 它们生成的主理想是单位理想.因此对每个 $j$ 存在类型为

\[
e_j + d_j = 1 \qquad \text{其中 } e_j \in \Big(\prod_{i \neq j} p_i^{n_i}\Big), \qquad d_j \in (p_j^{n_j})\]

的方程.于是如所希望, 对 $i \neq j$ 有 $\pi_i(e_j) = 0$, 对 $j$ 有 $\pi_j(e_j) = 1$.

$\varphi$ 的核由 $R$ 中被幂 $p_i^{n_i}$, $i = 1, \ldots, r$, 整除的一切元素组成, 从而由被 $a$ 整除的一切元素组成.于是 $\ker \varphi = (a)$, 由同态定理~\mtref{thm:5.1.10} 得到断言的同构.
\end{proof}

在定理~\mtref{thm:6.4.1} 的情形中, 剩余类环 $R/(a)$ 也可以看作 $R$-模.中国剩余定理的断言此时表明 $R/(a)$ 同构于 $R$-模 $R/(p_i^{n_i})$ 的构造性直和, 即

\[
R/(a) \simeq R/(p_1^{n_1}) \oplus \ldots \oplus R/(p_r^{n_r}),\]

其中如 \ref{sec:6.3} 节开头说明的那样, 把第 $i$ 个直和项 $R/(p_i^{n_i})$ 与相应的子模

\[
0 \times \ldots \times 0 \times R/(p_i^{n_i}) \times 0 \times \ldots \times 0 \subset R/(p_1^{n_1}) \times \ldots \times R/(p_r^{n_r})\]

等同.

对整环 $R$ 上的模 $M$, 通过

\[
T = \{a \in M;\; \text{存在 } \alpha \in R - \{0\} \text{ 使 } \alpha a = 0\}\]

定义所谓扭子模 $T$.利用 $R$ 的无零因子性容易验证 $T$ 确实是 $M$ 的子模.若 $M$ 与其扭子模一致, 则称 $M$ 为扭模.自由模的扭子模恒平凡, 因此自由模可以说应看作扭模的"对立面".

若 $M$ 是环 $R$ 上的模, 则 $R$-线性映射的正合序列

\[
R^n \xrightarrow{\psi} R^m \xrightarrow{\varphi} M \longrightarrow 0\]

其中 $R^m$ 与 $R^n$ 是合适的有限自由 $R$-模, 也称为 $M$ 的一个有限展示.更精确地说, 它由一个满态 $\varphi \colon R^m \longrightarrow M$ 和一个满足 $\operatorname{im} \psi = \ker \varphi$ 的 $R$-线性映射 $\psi \colon R^n \longrightarrow R^m$ 组成.当 $M$ 是主理想环 $R$ 上有限生成的模时, 这样的有限展示恒存在.只需在 $M$ 中选取一个有限生成系 $z_1, \ldots, z_m$, 并考虑把 $R^m$ 的典范基映到 $z_1, \ldots, z_m$ 的 $R$-线性映射 $\varphi \colon R^m \longrightarrow M$.则 $\varphi$ 是满态, 而按定理~\mtref{bigthm:6.3.4} 子模 $\ker \varphi \subset R^m$ 有限生成.因此可以找到另一个满态 $\psi \colon R^n \longrightarrow \ker \varphi \subset R^m$, 从而找到 $M$ 的一个有限展示.若在此情形中 $A \in R^{m \times n}$ 是关于 $R^m$ 与 $R^n$ 中合适基描述 $R$-线性映射 $\psi$ 的矩阵, 则我们也把 $A$ 的初等因子称为所考虑的 $M$ 的有限展示的初等因子.按推论~\mtref{cor:6.3.8}, 它们与 $\ker \varphi \subset R^m$ 的初等因子一致, 因此不依赖于 $R^m$ 与 $R^n$ 中基的选取.

我们现在想从初等因子定理~\mtref{bigthm:6.3.4} 推出关于主理想环上有限生成模的所谓\emph{主定理}(\textup{Hauptsatz})\index{zhudingli@主定理(Hauptsatz)}的各个版本.我们从初等因子定理的一个基本推论开始.

\begin{theorem}\label{thm:6.4.2}
设 $M$ 是主理想环 $R$ 上有限生成的模, $T \subset M$ 是其扭子模.则存在有限自由子模 $F \subset M$, 比如 $F \simeq R^d$, 以及非单位 $\alpha_1, \ldots, \alpha_s \in R - \{0\}$ 满足对 $1 \leq j < s$ 有 $\alpha_j \mid \alpha_{j + 1}$, 使得

\[
M = F \oplus T, \qquad T \simeq \bigoplus_{j = 1}^{s} R/\alpha_j R\text{.}\]

其中 $d$ 唯一确定; 称 $d$ 为 $M$ 的秩.此外, 元素 $\alpha_1, \ldots, \alpha_s$ 到差一个相伴为止唯一确定.它们与 $M$ 的每个有限展示的初等因子中的非单位一致.
\end{theorem}

\begin{proof}
我们从 $M$ 的一个有限展示出发, 考虑相应的满态 $\varphi \colon R^m \longrightarrow M$.则由同态定理有 $M \simeq R^m/\ker \varphi$.对情形 $\ker \varphi \subset R^m$ 可以应用初等因子定理~\mtref{bigthm:6.3.4}(其中考虑到按推论~\mtref{cor:6.3.7}, $R^m$ 的每个基恰好由 $m$ 个元素组成).因此存在构成 $R^m$ 的基的元素 $x_1, \ldots, x_m$ 以及元素 $\alpha_1, \ldots, \alpha_s \in R - \{0\}$, $s \leq m$, 满足对 $1 \leq j < s$ 有 $\alpha_j \mid \alpha_{j + 1}$, 使得 $\alpha_1 x_1, \ldots, \alpha_s x_s$ 构成 $\ker \varphi$ 的基.令 $\alpha_{s + 1} = \ldots = \alpha_m = 0$, 可以考虑满态

\[
\varphi' \colon R^m = \bigoplus_{j = 1}^{m} R x_j \longrightarrow \bigoplus_{j = 1}^{m} R x_j/R\alpha_j x_j, \quad (\gamma_1, \ldots, \gamma_m) \longmapsto (\bar{\gamma}_1, \ldots, \bar{\gamma}_m),\]

其中 $\bar{\gamma}_j$ 分别是 $\gamma_j$ 在 $Rx_j/R\alpha_j x_j$ 中的剩余类.按构造 $\ker \varphi' = \ker \varphi$, 从而由同态定理

\[
M \simeq R^m/\ker \varphi = \Big(\bigoplus_{j = 1}^{m} R x_j\Big)/\ker \varphi' \simeq \bigoplus_{j = 1}^{m} R x_j/R\alpha_j x_j \simeq R^{m - s} \oplus \bigoplus_{j = 1}^{s} R/\alpha_j R,\]

于是令 $d := m - s$ 就得到断言的分解, 只要忽略平凡项 $R/\alpha_j R = 0$, 即满足 $\alpha_j \in R^*$ 的项.

在这个分解中, $\bigoplus_{j = 1}^{s} R/\alpha_j R$ 对应于扭子模 $T \subset M$, 因此唯一确定.忽略平凡项 $R/\alpha_j R$, 即假设 $\alpha_1, \ldots, \alpha_s$ 不是单位, 则由引理~\mtref{lem:6.3.6} 得到 $\alpha_j$ 的唯一性.特别地, $\alpha_j$ 不依赖于所考虑的 $M$ 的展示, 并且按推论~\mtref{cor:6.3.8} 由这样的展示的初等因子中的非单位组成.

为看出 $d$ 也唯一确定, 考虑满态

\[
M \xrightarrow{\sim} R^d \oplus \bigoplus_{j = 1}^{s} R/\alpha_j R \longrightarrow R^d,\]

它由上面的同构与到直和项 $R^d$ 的投影复合而成.这个映射的核显然正好是扭子模 $T \subset M$, 因此由同态定理得到同构 $M/T \xrightarrow{\sim} R^d$.由此按推论~\mtref{cor:6.3.7} 得到 $d$ 的唯一性.
\end{proof}

关于主理想环上\emph{有限生成模的主定理}(\textup{Hauptsatz für endlich erzeugte Moduln über Hauptidealringen})\index{youxianshengchengmodezhudingli@有限生成模的主定理(Hauptsatz für endlich erzeugte Moduln über Hauptidealringen)}的更特殊版本, 可以借助中国剩余定理由定理~\mtref{thm:6.4.2} 推出.

\begin{theorem}\label{thm:6.4.3}
设 $M$ 是主理想环 $R$ 上有限生成的模, $T \subset M$ 是扭子模.再设 $P \subset R$ 是 $R$ 的素元素的一个代表系, 并对 $p \in P$, 用

\[
M_p = \{x \in M;\; p^n x = 0 \text{ \emph{对合}(\textup{Involution})\index{duihe@对合(Involution)}适的 } n \in \mathbb{N}\}\]

表示 $M$ 中所谓 $p$-\emph{挠子模}(\textup{Torsionsuntermodul})\index{naozimo@挠子模(Torsionsuntermodul)}.则

\[
T = \bigoplus_{p \in P} M_p,\]

并且存在有限自由子模 $F \subset M$, 比如 $F \simeq R^d$, 使

\[
M = F \oplus T,\]

其中 $d$ 唯一确定, 且对几乎一切 $p \in P$ 有 $M_p = 0$.此外, 对每个 $p \in P$ 存在自然数 $1 \leq n(p, 1) \leq \ldots \leq n(p, s_p)$ 使

\[
M_p \simeq \bigoplus_{j_p = 1 \ldots s_p} R/p^{n(p, j_p)} R,\]

其中 $s_p$, $n(p, j_p)$ 由同构

\[
M \simeq F \oplus \bigoplus_{p \in P} \bigoplus_{j_p = 1 \ldots s_p} R/p^{n(p, j_p)} R\]

唯一确定, 并且对几乎一切 $p$ 有 $s_p = 0$.
\end{theorem}

在进入证明之前, 我们还要把这个主定理专门表述为有限生成 $\mathbb{Z}$-模, 即关于\emph{有限生成阿贝尔群的主定理}(\textup{Hauptsatz für endlich erzeugte abelsche Gruppen})\index{youxianshengchengabeierqundezhudingli@有限生成阿贝尔群的主定理(Hauptsatz für endlich erzeugte abelsche Gruppen)}.

\begin{corollary}\label{cor:6.4.4}
设 $G$ 是有限生成的阿贝尔群, 即作为 $\mathbb{Z}$-模有限生成的阿贝尔群.设 $P$ 是 $\mathbb{N}$ 中素数的集合.则 $G$ 容许分解为子群

\[
G = F \oplus \bigoplus_{p \in P} \bigoplus_{j_p = 1 \ldots s_p} G_{p, j_p},\]

其中 $F$ 自由, 比如 $F \simeq \mathbb{Z}^d$, 并且 $G_{p, j_p} \simeq \mathbb{Z}/p^{n(p, j_p)}\mathbb{Z}$ 带有指数 $1 \leq n(p, 1) \leq \ldots \leq n(p, s_p)$; 其中 $G_{p, j_p}$ 称为阶为 $p$ 之幂的\emph{循环子群}(\textup{zyklische Gruppe; zyklische Untergruppe})\index{xunhuanziqun@循环子群(zyklische Gruppe; zyklische Untergruppe)}.数 $d$, $s_p$, $n(p, j_p)$ 由 $G$ 唯一确定, 子群 $G_p = \oplus_{j_p = 1 \ldots s_p} G_{p, j_p}$ 亦然, 其中对几乎一切 $p \in P$ 有 $s_p = 0$.
\end{corollary}

若 $G$ 是有限生成的\emph{挠}(\textup{Torsion})\index{nao@挠(Torsion)}群, 即 $\mathbb{Z}$ 上有限生成的扭模, 则 $G$ 没有自由部分, 从而如特别从推论~\mtref{cor:6.4.4} 所见只由有限多个元素组成.反之, 每个有限阿贝尔群自然是有限生成的挠群.

\begin{proof}[定理~\mtref{thm:6.4.3} 的证明]
我们从定理~\mtref{thm:6.4.2} 的分解

\[
M \simeq R^d \oplus \bigoplus_{j = 1}^{s} R/\alpha_j R\]

出发, 其中特别地 $d$ 唯一确定.这里 $\bigoplus_{j = 1}^{s} R/\alpha_j R$ 对应于扭子模 $T \subset M$, 而 $R^d$ 对应于自由模 $F \subset M$, 并且 $M = F \oplus T$.

把 $\alpha_j$ 分解为素因子, 比如

\[
\alpha_j = \varepsilon_j p_1^{n(1, j)} \ldots p_r^{n(r, j)}\]

其中 $\varepsilon_j \in R^*$ 是单位, $p_1, \ldots, p_r \in P$ 是两两不同的素元素, $n(i, j) \in \mathbb{N}$ 是指数, 也可以是平凡的.则由定理~\mtref{thm:6.4.1}

\[
T \simeq \bigoplus_{j = 1}^{s} \bigoplus_{i = 1}^{r} R/p_i^{n(i, j)} R \simeq \bigoplus_{i = 1}^{r} \bigoplus_{j = 1}^{s} R/p_i^{n(i, j)} R\text{.}\]

在这个分解中, $\bigoplus_{j = 1}^{s} R/p_i^{n(i, j)} R$ 显然正好对应于 $M$ 的 $p_i$-挠子模 $M_{p_i} \subset M$, 因此唯一确定; 因为 $p_i$ 的剩余类在每个形如 $R/p_{i'}^n R$(其中 $i' \neq i$) 的剩余类环中都是单位.于是由上面的分解特别得到 $T = \oplus_{p \in P} M_p$.此外, 在分解

\[
M_{p_i} \simeq \bigoplus_{j = 1}^{s} R/p_i^{n(i, j)} R\]

中, 对应指数 $n(i, j) = 0$ 的项是平凡的, 而其余项按引理~\mtref{lem:6.3.6} 唯一确定, 于是得到定理~\mtref{thm:6.4.3} 的全部断言.
\end{proof}

\aufgaben

对群 $G$, 用 ord $G$ 表示其元素个数, 并称之为 $G$ 的阶.

\noindent\textbf{习题 1.} 设 $G$ 是阶为 $n < \infty$ 的阿贝尔群.证明: 对 $n$ 的每个因子 $d$, 存在阶为 $d$ 的子群 $H \subset G$.另一方面, $G$ 的每个子群 $H \subset G$ 的阶都是 $n$ 的因子.(AT 354)

\medskip
\noindent\textbf{习题 2.} 对 $n \in \mathbb{N}$, 阶为 $n$ 的同构阿贝尔群的类至多有限多个.对 $n = 20$ 或 $n = 30$, 这个个数是多大?

\medskip
\noindent\textbf{习题 3.} 设 $a_1, \ldots, a_n \in \mathbb{N} - \{0\}$ 两两互素, $r_1, \ldots, r_n \in \mathbb{N}$ 是满足对 $i = 1, \ldots, n$ 有 $0 \leq r_i < a_i$ 的数.证明: 存在数 $a \in \mathbb{N}$, 使得 $a$ 被 $a_i$ 整数除时分别余 $r_i$.

\medskip
\noindent\textbf{习题 4.} 设 $M$ 是主理想环 $R$ 上有限生成的模.在 $M$ 不是扭模的情形, 证明: 存在有限多个自由子模 $F_1, \ldots, F_r \subset M$ 使 $M = \sum_{i = 1}^{r} F_i$.

\medskip
\noindent\textbf{习题 5.} 在主理想环 $R$ 上考虑带子模 $M \subset F$ 和 $M' \subset F'$ 的有限自由模 $F$ 和 $F'$, 使得剩余类模 $F/M$ 与 $F'/M'$ 都是扭模.证明: 存在同构 $F/M \xrightarrow{\sim} F'/M'$, 当且仅当 $M \subset F$ 的初等因子中的非单位与 $M' \subset F'$ 的初等因子中的非单位一致.

\medskip
\noindent\textbf{习题 6.} 设 $K$ 是域, $G \subset K^*$ 是 $K$ 的单位乘法群 $K^*$ 的有限子群.证明 $G$ 是循环的, 即存在 $n \in \mathbb{Z}$ 使 $G \simeq \mathbb{Z}/n\mathbb{Z}$.(AT 356)(提示: 考虑 $K$ 中类型为 $T^n - 1 \in K[T]$ 的多项式的零点.)


\section{矩阵的一般标准形与若尔当标准形}\index{biaozhunxing@标准形(Normalform)}\index{ruoerdangbiaozhunxing@若尔当标准形(Jordan'sche Normalform)}\label{sec:6.5}

设 $V$ 是域 $K$ 上的向量空间, $f \colon V \longrightarrow V$ 是自同态.正如 \ref{sec:6.3} 节所见, 在 $f$ 下可以把 $V$ 看作多项式环 $K[T]$ 上的模, 办法是把 $V$ 作为加法群,连同与 $K$ 中常数的纯量乘法, 但把与变量 $T$ 的乘法通过作用 $f$ 来定义.反过来, 可以从 $K[T]$-模 $M$ 重新得到带自同态 $f \colon V \longrightarrow V$ 的 $K$-向量空间 $V$.令 $V = M$, 并把 $M$ 的纯量乘法限制到常数多项式, 即限制到 $K$.这样 $V$ 成为 $K$-向量空间, 并且可以把 $K$-线性映射

\[
V \longrightarrow V, \qquad x \longmapsto T \cdot x,\]

看作 $V$ 的自同态.由于两个构造互为逆, 对给定的 $K$-向量空间 $V$ 得到 $K$-线性映射 $V \longrightarrow V$ 与 $V$ 上扩展现有 $K$-向量空间结构的 $K[T]$-模结构之间的双射.

现在把带自同态 $f \colon V \longrightarrow V$ 的有限维 $K$-向量空间 $V$ 看作 $K[T]$-模, 则存在唯一确定的次数最小的首一多项式 $p_f \in K[T]$ 使对一切 $v \in V$ 有 $p_f \cdot v = 0$; 称 $p_f$ 为 $f$ 的最小多项式, 参见定理~\mtref{thm:6.2.9}.此时 $V$ 特别是 $K[T]$-扭模, 甚至是有限生成的, 因为 $V$ 作为 $K$-向量空间有限生成.因此可以把结构定理~\mtref{thm:6.4.3} 应用于作为 $K[T]$-模的 $V$.

\begin{definition}\label{def:6.5.1}
设 $f \colon V \longrightarrow V$ 是 $K$-向量空间 $V$ 的自同态, $U \subset V$ 是线性子空间.
\begin{enumerate}
\item[(i)] 若 $f(U) \subset U$, 则称 $U$ 为 $f$-不变的.对由 $f$ 在 $U$ 上诱导的映射使用记号 $f|_U \colon U \longrightarrow U$;
\item[(ii)] 若存在向量 $u \in U$ 使序列 $u, f(u), f^2(u), \ldots$ 构成 $U$ 的生成系, 则称 $U$ 为 $f$-循环的; 特别地 $U$ 此时也是 $f$-不变的;
\item[(iii)] 若 $U$ 是 $f$-不变的线性子空间 $\neq 0$, 并且不能分解为两个真 $f$-不变的线性子空间的直和, 则称 $U$ 为 $f$-不可分解的.
\end{enumerate}
\end{definition}

我们现在想把刚才引入的子空间的性质也从模论上解释.环 $R$ 上的模 $M$ 称为单生成的, 如果它由一个元素生成, 即如果存在 $a \in M$ 使 $M = Ra$.

\begin{remark}\label{rem:6.5.2}
设 $f \colon V \longrightarrow V$ 是 $K$-向量空间 $V$ 的自同态, 其中把 $V$ 特别地也看作关于 $f$ 的 $K[T]$-模.则对子集 $U \subset V$ 有:
\begin{enumerate}
\item[(i)] $U$ 是 $V$ 的 $f$-不变的线性子空间, 当且仅当 $U$ 是 $V$ 的 $K[T]$-子模;
\item[(ii)] $U$ 是 $V$ 的 $f$-循环线性子空间, 当且仅当 $U$ 是 $V$ 的单生成 $K[T]$-子模;
\item[(iii)] $U$ 是 $V$ 的 $f$-不可分解线性子空间, 当且仅当 $U$ 是 $V$ 的 $K[T]$-子模 $\neq 0$, 并且不能分解为两个真 $K[T]$-子模的直和.
\end{enumerate}
\end{remark}

\begin{proof}
断言 (i) 直接清楚, 因为对线性子空间 $U \subset V$, 条件 $f(u) \in U$($u \in U$) 正好是 $T \cdot u \in U$, 因而等价于对 $q \in K[T]$, $u \in U$ 有 $q \cdot u \in U$.类似地得到断言 (ii), 因为对某个 $u \in U$, $U = \sum_{i \in \mathbb{N}} K \cdot f^i(u)$ 与 $U = \sum_{i \in \mathbb{N}} K \cdot T^i \cdot u$, 即 $U = K[T] \cdot u$ 等价.断言 (iii) 最后基于如下事实: 对 $f$-不变的线性子空间即 $K[T]$-子模 $U', U'' \subset V$, 和 $U = U' + U''$ 是直和, 当且仅当每个 $u \in U$ 有唯一确定的表示 $u = u' + u''$(其中 $u' \in U'$, $u'' \in U''$).至于在纯量乘法中允许取 $K$ 的元素还是 $K[T]$ 的元素, 是无关紧要的.
\end{proof}

在后面, 知道如何用 $K$-向量空间的同态来解释 $K[T]$-模的同态也很重要.

\begin{remark}\label{rem:6.5.3}
设 $V$ 与 $V'$ 是两个 $K$-向量空间, 它们关于给定的自同态 $f \colon V \longrightarrow V$ 与 $f' \colon V' \longrightarrow V'$ 看作 $K[T]$-模.对映射 $\varphi \colon V \longrightarrow V'$, 下列条件等价:
\begin{enumerate}
\item[(i)] $\varphi$ 是 $K[T]$-模的同态;
\item[(ii)] $\varphi$ 是 $K$-向量空间的同态, 并且下图

\begin{center}
\begin{tikzcd}[row sep=huge, column sep=huge]
V \arrow[r, "f"] \arrow[d, "\varphi"'] & V \arrow[d, "\varphi"] \\
V' \arrow[r, "f'"'] & V'
\end{tikzcd}
\end{center}

是交换的.
\end{enumerate}
相应的等价对单态,满态与同构 $\varphi \colon V \longrightarrow V'$ 也成立.特别地, 在同构 $\varphi$ 的情形, (i) 即 (ii) 中 $f$ 与 $f'$ 的最小多项式一致.$f$ 与 $f'$ 的特征多项式亦然.
\end{remark}

\begin{proof}
$K[T]$-模的每个同态 $\varphi \colon V \longrightarrow V'$ 是 $K[T]$-线性的, 从而特别是 $K$-线性的, 即 $K$-向量空间的同态.此外, (ii) 中图的交换性正好意味着 $\varphi$ 的 $T$-线性, 即对一切 $v \in V$ 有 $\varphi(T \cdot v) = T \cdot \varphi(v)$.于是由 (i) 得到 (ii); 反之, 若利用由 $T$-线性推出 $\varphi$ 对一切 $n \in \mathbb{N}$ 的 $T^n$-线性, 并且由可加性甚至推出 $\varphi$ 的 $K[T]$-线性, 则由 (ii) 也得到 (i).对单态,满态或同构的情形相应地推出.
\end{proof}

\begin{theorem}\label{thm:6.5.4}
设 $f \colon V \longrightarrow V$ 是维数 $r < \infty$ 的 $K$-向量空间 $V$ 的自同态.把 $V$ 看作关于 $f$ 的 $K[T]$-模.则下列条件等价:
\begin{enumerate}
\item[(i)] $V$ 是 $f$-循环的;
\item[(ii)] $V$ 作为 $K[T]$-模是单生成的;
\item[(iii)] 存在首一多项式 $p \in K[T]$, 使 $V$ 作为 $K[T]$-模同构于 $K[T]/(p)$.
\end{enumerate}
若这些条件成立, 则进一步有:
\begin{enumerate}
\item[(iv)] (iii) 中的多项式 $p$ 与 $f$ 的最小多项式 $p_f$ 及特征多项式 $\chi_f$ 一致; 特别地, $\operatorname{grad} p = r$.
\end{enumerate}
\end{theorem}

\begin{proof}
条件 (i) 与 (ii) 按注记~\mtref{rem:6.5.2} 等价.若 $V$ 再按条件 (ii) 作为 $K[T]$-模由元素 $u \in V$ 生成, 则考虑 $K[T]$-线性映射

\[
\varphi \colon K[T] \longrightarrow V, \quad h \longmapsto h \cdot u = h(f)(u);\]

$\varphi$ 是满射.此外, $\ker \varphi$ 是 $K[T]$ 的 $K[T]$-子模, 即理想, 由于 $K[T]$ 是主理想环, 它由某个多项式 $p \in K[T]$ 生成.由 $K[T]$-模的同态定理得到 (iii) 中希望的同构 $K[T]/(p) \xrightarrow{\sim} V$.由于 $\dim V < \infty$ 而 $\dim_K K[T] = \infty$, 有 $p \neq 0$, 并且可以把 $p$ 规定为首一多项式.

另一方面, 若条件 (iii) 成立, 则考虑剩余类映射 $K[T] \longrightarrow K[T]/(p)$, $h \longmapsto \bar{h}$.剩余类 $\bar{1} \in K[T]/(p)$ 自然作为 $K[T]$-模生成 $K[T]/(p)$, 即得到 (ii).

还剩 (iv) 需要论证, 其中我们把 (i) 即 (iii) 视为已成立.设 $\sigma \colon V \xrightarrow{\sim} K[T]/(p)$ 是按 (iii) 要求的 $K[T]$-模同构, 则由注记~\mtref{rem:6.5.3}, 下图

\begin{center}
\begin{tikzcd}[row sep=huge, column sep=huge]
V \arrow[r, "f"] \arrow[d, "\sigma"'] & V \arrow[d, "\sigma"] \\
K[T]/(p) \arrow[r, "f'"'] & K[T]/(p)
\end{tikzcd}
\end{center}

是交换的; $f'$ 表示 $K[T]/(p)$ 上与 $T$ 的乘法.我们想证明 $p$ 是 $f'$ 的最小多项式.按定理~\mtref{thm:6.2.9} 即定义~\mtref{def:6.2.10}, $f'$ 的最小多项式生成理想

\[
\begin{array}{l}
\big\{h \in K[T];\; h(f')(\bar{g}) = 0 \text{ 对一切 } \bar{g} \in K[T]/(p)\big\} \\
= \big\{h \in K[T];\; h \cdot \bar{g} = 0 \text{ 对一切 } \bar{g} \in K[T]/(p)\big\} \\
= \big\{h \in K[T];\; h \cdot \bar{1} = 0\big\} = \big\{h \in K[T];\; \bar{h} = 0\big\} \\
= (p) \subset K[T],
\end{array}\]

其中对 $h \in K[T]$, $\bar{h}$ 分别表示 $h$ 在 $K[T]/(p)$ 中的剩余类.故 $p$ 是 $f'$ 的最小多项式, 从而按注记~\mtref{rem:6.5.3} 与自同态 $f \colon V \longrightarrow V$ 的最小多项式 $p_f$ 一致.

$p = p_f$ 也可以用更详细的论证在自同态 $f$ 与向量空间 $V$ 的层面上看出: 按定理~\mtref{thm:6.2.9} 即定义~\mtref{def:6.2.10}, $p_f$ 生成 $K[T]$ 中由一切满足 $h(f) = 0$ 的多项式 $h$ 组成的理想.于是

\begin{equation}
(p_f) = \big\{h \in K[T];\; h(f)(v) = 0 \text{ 对一切 } v \in V\big\} \tag{$(*)$}\end{equation}

其中只需对 $V$ 的一个 $K$-生成系中的 $v$ 检验条件 $h(f)(v) = 0$.因此, 若 $V$ 作为 $K[T]$-模单生成, 比如由元素 $u \in V$ 生成, 则元素 $u, f(u), f^2(u), \ldots$ 构成 $V$ 的一个 $K$-生成系, 并且描述 $(*)$ 等价于

\begin{equation}
(p_f) = \big\{h \in K[T];\; h(f)\big(f^n(u)\big) = 0 \text{ 对一切 } n \in \mathbb{N}\big\} \text{.}\tag{$(**)$}\end{equation}

但计算

\[
h(f)\big(f^n(u)\big) = h \cdot f^n(u) = h \cdot T^n \cdot u = T^n \cdot h \cdot u = T^n \cdot h(f)(u)\]

表明, 从 $h(f)(u) = 0$ 已经推出对一切 $n \in \mathbb{N}$ 有 $h(f)(f^n(u)) = 0$.于是由 $(**)$ 得到

\[
(p_f) = \big\{h \in K[T];\; h(f)(u) = 0\big\} = \big\{h \in K[T];\; h \cdot u = 0\big\} = (p),\]

从而由于 $p_f$ 与 $p$ 的首一性已经有 $p_f = p$.

特别地, $p_f$ 的次数按下面的引理~\mtref{lem:6.5.5} 与 $K[T]/(p)$ 的 $K$-向量空间维数一致, 即与 $V$ 的维数 $r$ 一致.由于特征多项式 $\chi_f$ 是次数为 $\dim_K V$ 的首一多项式, 并且按\emph{凯莱--哈密顿定理}(\textup{Cayley-Hamilton, Satz von})\index{kailaihamidundingli@凯莱--哈密顿定理(Cayley-Hamilton, Satz von)} $p_f$ 是 $\chi_f$ 的因子(参见定理~\mtref{thm:6.2.11}), 已经有 $\chi_f = p_f$.不过方程 $\chi_f = p$ 也可以通过显式计算验证, 办法是把 $K[T]/(p)$ 上与 $T$ 的乘法关于基 $\bar{T}^0, \ldots, \bar{T}^{r-1}$ 用矩阵描述并求其特征多项式.
\end{proof}

最后还要计算类型为 $K[T]/(p)$ 的商的 $K$-向量空间维数.

\begin{lemma}\label{lem:6.5.5}
设 $p \in K[T]$ 是首一多项式.则

\[
\dim_K\big(K[T]/(p)\big) = \operatorname{grad} p,\]

并且剩余类 $\bar{T}^0, \ldots, \bar{T}^{s - 1}$(其中 $s = \operatorname{grad} p$) 构成 $K[T]/(p)$ 的一个 $K$-基.
\end{lemma}

\begin{proof}
只需证明对 $s = \operatorname{grad} p$, 剩余类 $\bar{T}^0, \ldots, \bar{T}^{s - 1}$ 构成 $K[T]/(p)$ 的一个 $K$-基.为此考虑任意元素 $\bar{h} \in K[T]/(p)$ 并选取一个原像 $h \in K[T]$.带余除法给出分解 $h = qp + c$, 其中 $q, c \in K[T]$ 是多项式, $c$ 的次数 $< s$, 比如 $c = \sum_{i = 0}^{s - 1} c_i T^i$ 且系数 $c_i \in K$.于是

\[
\bar{h} = \bar{c} = \sum_{i = 0}^{s - 1} c_i \bar{T}^i,\]

即 $K[T]/(p)$ 作为 $K$-向量空间由剩余类 $\bar{T}^0, \ldots, \bar{T}^{s - 1}$ 生成.但这些元素在 $K$ 上也是线性无关的.因为若有 $\sum_{i = 0}^{s - 1} c_i \bar{T}^i = 0$ 且系数 $c_i \in K$, 则这意味着 $\sum_{i = 0}^{s - 1} c_i T^i \in (p)$, 从而对一切 $i = 0, \ldots, s - 1$ 有 $c_i = 0$, 因为 $(p)$ 中一切非零元素都有次数 $\operatorname{grad} p = s$.
\end{proof}

我们现在想把结构定理~\mtref{thm:6.4.3} 改写为 $K[T]$-模的结构定理, 或者更确切地说, 改写为带自同态 $f \colon V \longrightarrow V$ 的有限维 $K$-向量空间 $V$ 情形的结构定理.

\begin{bigtheorem}\label{bigthm:6.5.6}
设 $V$ 是带自同态 $f \colon V \longrightarrow V$ 的有限维 $K$-向量空间.则存在两两不同的首一(从而不相伴的)素多项式 $p_1, \ldots, p_r \in K[T]$, 自然数 $1 \leq n(i, 1) \leq \ldots \leq n(i, s_i)$, 以及 $f$-循环的 $K$-线性子空间 $V_{ij} \subset V$, $i = 1, \ldots, r$, $j = 1, \ldots, s_i$, 使得在 $K[T]$-模的意义下

\[
V = \bigoplus_{i = 1}^{r} \bigoplus_{j = 1}^{s_i} V_{ij}, \qquad V_{ij} \simeq K[T]/(p_i^{n(i, j)})\text{.}\]

特别地, $p_i^{n(i, j)}$ 是 $f|_{V_{ij}}$ 的最小多项式即特征多项式; 参见定理~\mtref{thm:6.5.4}.

在上述分解中, 素多项式 $p_1, \ldots, p_r$, 数 $n(i, j)$ 以及向量空间 $V_i := \oplus_{j = 1}^{s_i} V_{ij}$, $i = 1, \ldots, r$, 唯一确定, 但子空间 $V_{ij}$ 本身未必唯一; 有

\[
V_i = \bigcup_{n \in \mathbb{N}} \ker p_i^n(f)\text{.}\]

此外, $p_f = p_1^{n(1, s_1)} \cdot \ldots \cdot p_r^{n(r, s_r)}$ 是 $f$ 的最小多项式, 并且

\[
\dim_K V_{ij} = n(i, j) \cdot \operatorname{grad} p_i,\]

特别地

\[
\dim_K V = \sum_{i = 1}^{r} \sum_{j = 1}^{s_i} n(i, j) \cdot \operatorname{grad} p_i\text{.}\]
\end{bigtheorem}

\begin{proof}
把 $V$ 分解为所述类型的 $f$-\emph{循环子空间}(\textup{zyklischer Untervektorraum})\index{xunhuanzikongjian@循环子空间(zyklischer Untervektorraum)} $V_{ij}$ 的存在性与唯一性断言直接由定理~\mtref{thm:6.4.3} 和定理~\mtref{thm:6.5.4} 得出.其中 $V_i = \bigoplus_{j = 1}^{s_i} V_{ij}$ 用 $K[T]$-模的语言说就是 $p_i$-挠子模, 而定理~\mtref{thm:6.4.3} 中的自由部分消失, 因为 $V$ 是 $K[T]$-扭模.

按定理~\mtref{thm:6.5.4}, $p_i^{n(i, j)}$ 是 $f$ 在 $V_{ij}$ 上限制的最小多项式.则容易看出, 所有这些多项式的最小公倍式, 即

\[
p_1^{n(1, s_1)} \cdot \ldots \cdot p_r^{n(r, s_r)}\]

是 $f$ 在 $V$ 上的最小多项式.维数公式最后利用推论~\mtref{cor:1.6.6} 由引理~\mtref{lem:6.5.5} 得出.
\end{proof}

\begin{corollary}\label{cor:6.5.7}
设 $f \colon V \longrightarrow V$ 是有限维 $K$-向量空间 $V$ 的自同态.则下列条件等价:
\begin{enumerate}
\item[(i)] $V$ 是 $f$-不可分解的;
\item[(ii)] $V$ 是 $f$-循环的, 并且 $p_f = \chi_f$ 是 $K[T]$ 中某个素多项式的幂.
\end{enumerate}
特别地, 定理~\mtref{bigthm:6.5.6} 中的 $f$-循环 $K$-线性子空间 $V_{ij} \subset V$ 是 $f$-不可分解的.
\end{corollary}

\begin{proof}
按定理~\mtref{bigthm:6.5.6} 把 $V$ 分解为 $f$-循环子空间的直和, 其最小多项式分别是 $K[T]$ 中素多项式的幂.若 $V$ 是 $f$-不可分解的, 则这个分解只能由一个直和项组成, 于是推出 $V$ 是 $f$-循环的.

另一方面, 若 $V$ 是 $f$-循环的, 且最小多项式 $p_f$ 是 $K[T]$ 中素多项式的幂, 则按定理~\mtref{thm:6.5.4} 在 $K[T]$-模的意义下 $V \simeq K[T]/(p_f)$, 而定理~\mtref{bigthm:6.5.6} 中分解的唯一性表明 $V$ 是 $f$-不可分解的.
\end{proof}

最后我们想把分解定理~\mtref{bigthm:6.5.6} 改写为关于方阵的标准形的断言.我们从一条平凡的观察开始, 其中使用下面的记号: 对矩阵 $A_i \in K^{n_i \times n_i}$, $i = 1, \ldots, r$, 用 $\operatorname{Diag}(A_1, \ldots, A_r) \in K^{n \times n}$, $n = \sum_{i = 1}^{r} n_i$, 表示以 $A_1, \ldots, A_r$ 为"对角元素"的"对角矩阵", 即

\[
\operatorname{Diag}(A_1, \ldots, A_r) = \left(\begin{array}{ccccc}
A_1 & & & & 0 \\
 & A_2 & & & \\
 & & \ldots & & \\
 & & & \ldots & \\
0 & & & & A_r
\end{array}\right)\]

若对 $i = 1, \ldots, r$ 有 $A_i = (\lambda_i) \in K^{1 \times 1}$, 则 $\operatorname{Diag}(A_1, \ldots, A_r)$ 是真正的对角矩阵, 并且我们用 $\operatorname{Diag}(\lambda_1, \ldots, \lambda_r)$ 代替 $\operatorname{Diag}(A_1, \ldots, A_r)$.

\begin{lemma}\label{lem:6.5.8}
设 $f \colon V \longrightarrow V$ 是有限维 $K$-向量空间 $V$ 的自同态, 带分解 $V = \bigoplus_{i = 1}^{r} V_i$ 为 $f$-不变的子空间.设 $X_i$ 分别是 $V_i$ 的基, 而 $X$ 是由基 $X_1, \ldots, X_r$ 合成的 $V$ 的基.令 $f_i = f|_{V_i}$, 则对属于 $f$ 的矩阵以及特征多项式即最小多项式有:
\begin{enumerate}
\item[(i)] $A_{f, X, X} = \operatorname{Diag}(A_{f_1, X_1, X_1}, \ldots, A_{f_r, X_r, X_r})$;
\item[(ii)] $\chi_f = \chi_{f_1} \cdot \ldots \cdot \chi_{f_r}$;
\item[(iii)] $p_f = \mathrm{kgV}(p_{f_1}, \ldots, p_{f_r})$.
\end{enumerate}
\end{lemma}

\begin{lemma}\label{lem:6.5.9}
设 $f \colon V \longrightarrow V$ 是带基 $X = (x_1, \ldots, x_n)$ 的 $K$-向量空间 $V$ 的自同态.若 $A_{f, X, X} = \operatorname{Diag}(A_1, \ldots, A_r)$ 带矩阵 $A_i \in K^{n_i \times n_i}$(其中 $n = \sum_{i = 1}^{r} n_i$), 则对 $m_i = \sum_{j = 1}^{i} n_j$, 子空间

\[
V_i = \sum_{j = m_{i - 1} + 1}^{m_i} K \cdot x_j, \quad i = 1, \ldots, r,\]

是 $f$-不变的, 并且 $V = \bigoplus_{i = 1}^{r} V_i$.
\end{lemma}

\begin{proof}[引理~\mtref{lem:6.5.8} 与引理~\mtref{lem:6.5.9} 的证明]
引理~\mtref{lem:6.5.8} 中的断言 (i) 以及引理~\mtref{lem:6.5.9} 由定义~\mtref{def:3.1.3} 中定义的线性映射与描述矩阵之间的对应直接清楚.此外, 引理~\mtref{lem:6.5.8} 中的断言 (ii) 由 \ref{sec:4.3} 节末尾的例 (2)(参见例~\mtref{bsp:4.3.2}) 得出.断言 (iii) 最后成立, 因为一方面所有 $p_{f_i}$ 的最小公倍式消去自同态 $f$, 另一方面 $p_f$ 必须是每个 $p_{f_i}$ 的倍式.
\end{proof}

对次数为 $n$ 的首一多项式

\[
p = T^n + \sum_{i = 0}^{n - 1} c_i T^i \in K[T],\]

称矩阵

\[
A(p) = \left(\begin{array}{cccccc}
0 & & & & & -c_0 \\
1 & 0 & & & & -c_1 \\
 & 1 & \ddots & & & -c_2 \\
 & & \ddots & \ddots & & \vdots \\
 & & & \ddots & 0 & \vdots \\
 & & & & 1 & -c_{n-1}
\end{array}\right)\]

为 $p$ 的友矩阵.特别地, 对 $n = 1$ 有 $A(p) = (-c_0)$.我们也不想排除 $n = 0$ 的特例, 此时 $A(p)$ 是空矩阵.

\begin{lemma}\label{lem:6.5.10}
设 $f \colon V \longrightarrow V$ 是维数 $n < \infty$ 的 $K$-向量空间的自同态, $p \in K[T]$ 是次数为 $n$ 的首一多项式.则下列条件等价:
\begin{enumerate}
\item[(i)] $V$ 是 $f$-循环的, 且最小多项式 $p_f = p$.(注意, 若 $\operatorname{grad} p_f = \dim_K V$, 即 $p_f = \chi_f$, 则 $V$ 自动是 $f$-循环的.这一点稍后将由定理~\mtref{bigthm:6.5.18} 读出.)
\item[(ii)] 存在 $V$ 的基 $X$, 使矩阵 $A_{f, X, X}$ 是 $p$ 的友矩阵.
\end{enumerate}
\end{lemma}

\begin{proof}
对 $V = 0$, 断言平凡, 所以可设 $V \neq 0$.首先设 $V$ 是 $f$-循环的, 即在 $K[T]$-模的意义下 $V \simeq K[T]/(p)$, 其中 $p = p_f$ 是 $f$ 的最小多项式; 参见定理~\mtref{thm:6.5.4}.则按引理~\mtref{lem:6.5.5}, 剩余类 $\bar{T}^0, \ldots, \bar{T}^{n - 1}$ 构成 $K[T]/(p)$ 的一个 $K$-基; 换言之, 存在元素 $u \in V$(即同构 $V \simeq K[T]/(p)$ 下与剩余类 $\bar{1} \in K[T]/(p)$ 对应的元素), 使得 $f^0(u), f^1(u), \ldots, f^{n - 1}(u)$ 构成 $V$ 的一个 $K$-基 $X$.对 $p = T^n + c_{n - 1} T^{n - 1} + \ldots + c_0$ 于是

\[
f\big(f^{n - 1}(u)\big) = f^n(u) = -\sum_{i = 0}^{n - 1} c_i f^i(u),\]

并且看到 $A_{f, X, X}$ 是 $p$ 的友矩阵.

反之, 设 $X = (x_1, \ldots, x_n)$ 是 $V$ 的基, 使得 $A_{f, X, X}$ 是某个次数为 $n$ 的首一多项式 $p \in K[T]$ 的友矩阵.令 $u = x_1$, 则对 $i = 1, \ldots, n$ 有 $x_i = f^{i - 1}(u)$, 于是推出 $V$ 是 $f$-循环的.此外, 由矩阵 $A_{f, X, X}$ 读出 $p(f)(u) = 0$, 即在 $V$ 作为 $K[T]$-模的意义下 $p \cdot u = 0$.但由于 $V = K[T] \cdot u$, 也得到 $p \cdot V = 0$, 所以最小多项式 $p_f$ 是 $p$ 的因子.由于按定理~\mtref{thm:6.5.4}, $p_f$ 与特征多项式 $\chi_f$ 一致, 因而与 $p$ 一样次数为 $n$, 于是已经有 $\chi_f = p_f = p$.
\end{proof}

现在我们可以从定理~\mtref{bigthm:6.5.6} 推导所谓矩阵的一般标准形的存在性与唯一性.

\begin{bigtheorem}[一般标准形]\label{bigthm:6.5.11}\index{yibanbiaozhunxing@一般标准形(allgemeine Normalform)}
每个矩阵 $A \in K^{n \times n}$ 都相似于形如 $\operatorname{Diag}(A_1, \ldots, A_r)$ 的矩阵, 其中 $A_i$ 是方阵, 并且 $A_i$ 分别是 $K[T]$ 中某个首一素多项式之幂 $q_i$ 的友矩阵.矩阵 $A_1, \ldots, A_r$ 到差一个次序为止由 $A$ 唯一确定.对 $A$ 的最小多项式有 $p_A = \mathrm{kgV}(q_1, \ldots, q_r)$, 对 $A$ 的特征多项式有 $\chi_A = q_1 \cdot \ldots \cdot q_r$.
\end{bigtheorem}

\begin{proof}
把 $A$ 解释为有限维 $K$-向量空间 $V$ 的自同态 $f \colon V \longrightarrow V$ 的矩阵, 则按定理~\mtref{bigthm:6.5.6}, $V$ 分解为 $f$-循环子空间 $V_1, \ldots, V_r$ 的直和, 比如 $V_i \simeq K[T]/(q_i)$, 其中 $q_i$ 是某个首一素多项式的幂.按引理~\mtref{lem:6.5.8} 这表明 $A$ 相似于对角矩阵 $\operatorname{Diag}(A_1, \ldots, A_r)$, 其中 $A_i$ 分别关于 $V_i$ 的某个合适基描述自同态 $f_i = f|_{V_i}$.由定理~\mtref{thm:6.5.4} 得出 $q_i$ 是 $f_i$ 的最小多项式即特征多项式, 而由引理~\mtref{lem:6.5.10} 得出可以把 $A_i$ 取作 $q_i$ 的友矩阵.

反之, 若 $A$ 相似于形如 $\operatorname{Diag}(A_1, \ldots, A_r)$ 的矩阵, 其中 $A_i$ 分别是某个首一素多项式之幂 $q_i$ 的友矩阵, 则它对应于分解 $V \simeq \oplus_{i = 1}^{r} K[T]/(q_i)$; 参见引理~\mtref{lem:6.5.9} 和引理~\mtref{lem:6.5.10} 连同定理~\mtref{thm:6.5.4}.定理~\mtref{bigthm:6.5.6} 中的唯一性断言于是蕴含矩阵 $A_1, \ldots, A_r$ 的唯一性, 因为它们由它们的最小多项式即特征多项式唯一确定.

关于最小多项式 $p_A$ 与特征多项式 $\chi_A$ 的断言已经在引理~\mtref{lem:6.5.8} 中推导过.
\end{proof}

作为推论可以再次读出凯莱--哈密顿定理, 甚至是在一个稍加改进的版本中:

\begin{corollary}[凯莱--哈密顿]\label{cor:6.5.12}
设 $f \colon V \longrightarrow V$ 是维数 $n < \infty$ 的 $K$-向量空间的自同态, 最小多项式为 $p_f$, 特征多项式为 $\chi_f$.则 $p_f$ 是 $\chi_f$ 的因子, 并且 $\chi_f$ 的素因子分解中出现的每个素因子也在 $p_f$ 的素因子分解中出现, 不过一般以较低的幂次.
\end{corollary}

对角矩阵 $A = \operatorname{Diag}(\lambda_1, \ldots, \lambda_n) \in K^{n \times n}$ 已经是一般标准形, 因为矩阵 $A_i = (\lambda_i) \in K^{1 \times 1}$ 可以看作多项式 $T - \lambda_i \in K[T]$ 的友矩阵.由定理~\mtref{bigthm:6.5.11} 因此得到如下的可对角化判别准则:

\begin{corollary}\label{cor:6.5.13}
对有限维 $K$-向量空间的自同态 $f$(或矩阵 $A \in K^{n \times n}$), 下列断言等价:
\begin{enumerate}
\item[(i)] $f$(即 $A$) 可对角化;
\item[(ii)] 最小多项式 $p_f$(即 $p_A$) 分解为乘积 $\prod_{i = 1}^{r}(T - \lambda_i)$, 其中 $\lambda_1, \ldots, \lambda_r \in K$ 是两两不同的零点.
\end{enumerate}
\end{corollary}


接下来要推导矩阵的所谓\emph{若尔当标准形}(\textup{Jordan'sche Normalform})\index{ruoerdangbiaozhunxing@若尔当标准形(Jordan'sche Normalform)}.这种标准形只能在特征多项式即最小多项式完全分解为线性因子的情形下构造, 例如当域 $K$ 代数闭时.对 $\lambda \in K$, 称矩阵

\[
J(\lambda, n) = \left(\begin{array}{ccccccc}
\lambda & & & & & & 0 \\
1 & \lambda & & & & & \\
 & 1 & \lambda & & & & \\
 & & \cdot & \cdot & & & \\
 & & & \cdot & \cdot & & \\
 & & & & 1 & \lambda & \\
0 & & & & & 1 & \lambda
\end{array}\right) \in K^{n \times n}\]

为属于特征值 $\lambda$ 的长度为 $n$ 的若尔当块.

\begin{lemma}\label{lem:6.5.14}
对 $A \in K^{n \times n}$ 和 $\lambda \in K$, 下列条件等价:
\begin{enumerate}
\item[(i)] $A$ 相似于 $J(\lambda, n)$;
\item[(ii)] $A$ 相似于多项式 $q = (T - \lambda)^n \in K[T]$ 的友矩阵 $A(q)$.
\end{enumerate}
\end{lemma}

\begin{proof}
把 $A$ 解释为自同态 $f \colon V \longrightarrow V$ 的描述矩阵, 其中 $V$ 是 $n$ 维 $K$-向量空间.若 $A$ 相似于 $J(\lambda, n)$, 则存在 $V$ 的 $K$-基 $x_1, \ldots, x_n$ 使

\begin{equation}
f(x_i) = \lambda x_i + x_{i + 1}, \qquad i = 1, \ldots, n, \tag{$(*)$}\end{equation}

其中令 $x_{n + 1} = 0$.由此用归纳法得到

\[
x_1, \ldots, x_i \in \big\langle f^0(x_1), \ldots, f^{i - 1}(x_1)\big\rangle, \qquad i = 1, \ldots, n;\]

因为对 $i = 1$ 这是平凡的, 而利用对指标 $i$ 的归纳假设推出

\[
\begin{array}{l}
x_{i + 1} = f(x_i) - \lambda x_i \\
\qquad \in \big\langle f^1(x_1), \ldots, f^i(x_1)\big\rangle + \big\langle f^0(x_1), \ldots, f^{i - 1}(x_1)\big\rangle \\
\qquad = \big\langle f^0(x_1), \ldots, f^i(x_1)\big\rangle\text{.}
\end{array}\]

特别地因此

\[
V = \big\langle f^0(x_1), \ldots, f^{n - 1}(x_1)\big\rangle,\]

即 $V$ 是 $f$-循环的, 并作为 $K[T]$-模由 $x_1$ 生成.现在考虑满射的 $K[T]$-模同态

\[
\varphi \colon K[T] \longrightarrow V, \qquad h \longmapsto h \cdot x_1 = h(f)(x_1)\text{.}\]

由 $(*)$ 得到

\[
(f - \lambda \operatorname{id})(x_i) = x_{i + 1}, \qquad i = 1, \ldots, n,\]

即

\[
(f - \lambda \operatorname{id})^i(x_1) = \left\{\begin{array}{l l} x_{i + 1} \neq 0 & \text{对 } i = 0, \ldots, n - 1, \\ 0 & \text{对 } i = n \text{.} \end{array}\right.\]

因此 $\varphi$ 的核含有多项式 $(T - \lambda)^n$, 但不含 $(T - \lambda)^{n - 1}$.作为主理想, $\ker \varphi$ 由某个首一多项式 $q \in K[T]$ 生成.它整除 $(T - \lambda)^n$, 但不整除 $(T - \lambda)^{n - 1}$, 因而与 $(T - \lambda)^n$ 一致.由模的同态定理, $\varphi$ 于是在 $K[T]$-模的意义下诱导同构 $K[T]/((T - \lambda)^n) \xrightarrow{\sim} V$, 并且按定理~\mtref{thm:6.5.4} 把 $V$ 视为最小多项式 $p_f = (T - \lambda)^n$ 的 $f$-循环模.最后按引理~\mtref{lem:6.5.10}, $A$ 相似于最小多项式 $(T - \lambda)^n$ 的友矩阵.因此蕴含 (i) $\Longrightarrow$ (ii) 得证.

为证明反向蕴含, 设 $A$ 相似于多项式 $q = (T - \lambda)^n \in K[T]$ 的友矩阵.则按引理~\mtref{lem:6.5.10}, $V$ 是 $f$-循环的, 且最小多项式为 $q$, 并且在 $K[T]$-模的意义下 $V \simeq K[T]/(q)$.我们想证明剩余类

\[
(\bar{T} - \lambda)^0, (\bar{T} - \lambda)^1, \ldots, (\bar{T} - \lambda)^{n - 1}\]

构成 $K[T]/(q)$ 的一个 $K$-基.按引理~\mtref{lem:6.5.5} 已知这对幂 $\bar{T}^0, \ldots, \bar{T}^{n - 1}$ 成立.因此只需归纳证明

\[
\bar{T}^0, \ldots, \bar{T}^{i - 1} \in \big\langle (\bar{T} - \lambda)^0, \ldots, (\bar{T} - \lambda)^{i - 1}\big\rangle, \qquad i = 1, \ldots, n\text{.}\]

对 $i = 1$ 这是平凡的, 而利用对指标 $i - 1$ 的归纳假设可以如下推出:

\[
\begin{array}{l}
\bar{T}^i = (\bar{T} - \lambda) \cdot \bar{T}^{i - 1} + \lambda \bar{T}^{i - 1} \\
\quad \in \big\langle (\bar{T} - \lambda)^1, \ldots, (\bar{T} - \lambda)^i\big\rangle + \big\langle (\bar{T} - \lambda)^0, \ldots, (\bar{T} - \lambda)^{i - 1}\big\rangle \\
\quad = \big\langle (\bar{T} - \lambda)^0, \ldots, (\bar{T} - \lambda)^i\big\rangle
\end{array}\]

因此幂 $(\bar{T} - \lambda)^{i - 1}$, $i = 1, \ldots, n$, 构成 $K[T]/(q)$ 的一个 $K$-基.

现在考虑方程

\[
\bar{T} \cdot (\bar{T} - \lambda)^{i - 1} = \lambda \cdot (\bar{T} - \lambda)^{i - 1} + (\bar{T} - \lambda)^i, \qquad i = 1, \ldots, n,\]

其中 $(\bar{T} - \lambda)^n = 0$.它们表明, 按注记~\mtref{rem:6.5.3} 与 $f$ 对应的自同态

\[
K[T]/(q) \longrightarrow K[T]/(q), \quad \bar{h} \longmapsto \bar{T} \cdot \bar{h},\]

关于上述基由若尔当块 $J(\lambda, n)$ 描述.于是 $V$ 中也存在一个基, 关于它 $f$ 由 $J(\lambda, n)$ 描述, 因而蕴含 (ii) $\Longrightarrow$ (i) 得证.
\end{proof}

现在可以容易地从定理~\mtref{bigthm:6.5.11} 推出若尔当标准形的存在性与唯一性.

\begin{bigtheorem}[若尔当标准形]\label{bigthm:6.5.15}\index{ruoerdangbiaozhunxing@若尔当标准形(Jordan'sche Normalform)}
设 $A \in K^{n \times n}$ 是最小多项式即特征多项式完全分解为线性因子的矩阵.则 $A$ 相似于所谓若尔当矩阵

\[
\operatorname{Diag}\big(J(\lambda_1, n_1), \ldots, J(\lambda_r, n_r)\big),\]

其"对角元素"是若尔当块.元素 $\lambda_i, n_i$ 到差一个次序为止由 $A$ 唯一确定, 其中 $\lambda_i$(在可能的重复列举下) 恰好遍历 $A$ 的特征值.若选取 $n$ 维 $K$-向量空间 $V$ 并把 $A$ 实现为自同态 $f \colon V \longrightarrow V$ 的矩阵, 则

\[
V \simeq \bigoplus_{i = 1}^{r} K[T]/\big((T - \lambda_i)^{n_i}\big)\]

正好是定理~\mtref{thm:6.4.3} 即定理~\mtref{bigthm:6.5.6} 中的分解.
\end{bigtheorem}

\begin{proof}
利用引理~\mtref{lem:6.5.14}, 从定理~\mtref{bigthm:6.5.11} 中 $A$ 的一般标准形可以到达若尔当标准形, 反之从若尔当标准形也可以到达一般标准形.
\end{proof}

我们说矩阵 $A \in K^{n \times n}$ \emph{可三角化}(\textup{trigonalisierbar})\index{kesanjiaohua@可三角化(trigonalisierbar)}, 如果它相似于形如 $B = (\beta_{ij}) \in K^{n \times n}$ 且对 $i < j$ 有 $\beta_{ij} = 0$ 的矩阵; 称 $B$ 为(下)三角矩阵.由若尔当标准形的存在性可以特别推出三角化的判别准则.

\begin{corollary}\label{cor:6.5.16}
矩阵 $A \in K^{n \times n}$ 可三角化, 当且仅当 $A$ 的最小多项式即特征多项式完全分解为线性因子.
\end{corollary}

\begin{proof}
设 $A$ 可三角化, 即相似于对 $i < j$ 有 $\beta_{ij} = 0$ 的矩阵 $B = (\beta_{ij}) \in K^{n \times n}$.则 $\chi_A = \chi_B = \prod_{i = 1}^{n}(T - \beta_{ii})$; 特别地 $\chi_A$ 从而按推论~\mtref{cor:6.5.12} 也最小多项式 $p_A$ 完全分解为线性因子.反之, 若后一条件成立, 则由若尔当标准形的存在性看出 $A$ 可三角化.
\end{proof}

作为下一步, 我们给出一个显式计算矩阵的若尔当标准形的第一个(相当粗略的)实用方法.考虑特征多项式完全分解为线性因子的矩阵 $A \in K^{n \times n}$, 比如

\[
\chi_A = (T - \lambda_1)^{n_1} \cdot \ldots \cdot (T - \lambda_r)^{n_r};\]

与定理~\mtref{bigthm:6.5.15} 不同, 这里假设特征值 $\lambda_1, \ldots, \lambda_r$ 两两不同.$A$ 的若尔当标准形 $J$ 是对角线上排列着类型为 $J(\lambda, s)$ 的若尔当块的矩阵.对 $i = 1, \ldots, r$ 和 $j = 1, \ldots, n_i$, 用 $k_{i, j}$ 表示若尔当块 $J(\lambda_i, j)$ 在 $J$ 中出现的次数.有

\[
k_{i, 1} + 2 k_{i, 2} + \ldots + n_i k_{i, n_i} = n_i,\]

因为特征值 $\lambda_i$ 的重数正好是 $n_i$, 即元素 $\lambda_i$ 在 $J$ 的对角线上恰好出现 $n_i$ 次.

由于矩阵 $A$ 与 $J$ 相似, 矩阵 $A - \lambda E$ 与 $J - \lambda E$ 以及幂 $(A - \lambda E)^\ell$ 与 $(J - \lambda E)^\ell$($\ell \in \mathbb{N}$) 也相似; 其中 $\lambda \in K$, $E \in K^{n \times n}$ 为单位矩阵.特别地, 按定理~\mtref{thm:3.4.6} 有

\[
\operatorname{rg}(A - \lambda E)^\ell = \operatorname{rg}(J - \lambda E)^\ell \quad \text{对 } \ell \in \mathbb{N},\]

而 $(J - \lambda E)^\ell$ 又是由方块, 即 $J - \lambda E$ 的若尔当块的 $\ell$ 次幂组成的"对角矩阵".现在对类型为 $J(\lambda, m)$ 的若尔当块, $\ell$ 次幂的秩算得为

\[
\operatorname{rg} J(\lambda, m)^\ell = \left\{\begin{array}{l l} \max\{0, m - \ell\} & \text{对 } \lambda = 0, \\ m & \text{对 } \lambda \neq 0 \text{.} \end{array}\right.\]

因此成立如下的方程:

\[
\begin{array}{r l}
& \operatorname{rg}(A - \lambda_i E)^{n_i} = \operatorname{rg}(J - \lambda_i E)^{n_i} = n - n_i, \\
& \operatorname{rg}(A - \lambda_i E)^{n_i - 1} = \operatorname{rg}(J - \lambda_i E)^{n_i - 1} = n - n_i + k_{i, n_i}, \\
& \operatorname{rg}(A - \lambda_i E)^{n_i - 2} = \operatorname{rg}(J - \lambda_i E)^{n_i - 2} = n - n_i + 2 k_{i, n_i} + k_{i, n_i - 1},
\end{array}\]

\[
\operatorname{rg}(A - \lambda_i E)^1 = \operatorname{rg}(J - \lambda_i E)^1 = n - n_i + (n_i - 1) k_{i, n_i} + \ldots + k_{i, 2}\text{.}\]

因此, 若能确定 $A - \lambda_i E$ 的幂的秩, 就可以依次计算数 $k_{i, j}$, $j = n_i, n_i - 1, \ldots, 2$.此外, 由上面的方程还得到

\[
k_{i, 1} = n_i - 2 k_{i, 2} - \ldots - n_i k_{i, n_i},\]

所以也可以计算 $k_{i, 1}$.但一旦所有数 $k_{i, j}$ 已知, 就完整地确定了 $A$ 的若尔当标准形 $J$.

作为例子, 我们想确定矩阵

\[
A = \left(\begin{array}{cccc} 1 & 1 & 0 & 1 \\ 0 & 2 & 0 & 0 \\ -1 & 1 & 2 & 1 \\ -1 & 1 & 0 & 3 \end{array}\right) \in \mathbb{R}^{4 \times 4}\]

的若尔当标准形.在上述记号中有 $\chi_A = (T - 2)^4$, 即 $r = 1$, $n_1 = 4$.我们有

\[
\operatorname{rg}(A - 2E) = \operatorname{rg}\left(\begin{array}{cccc} -1 & 1 & 0 & 1 \\ 0 & 0 & 0 & 0 \\ -1 & 1 & 0 & 1 \\ -1 & 1 & 0 & 1 \end{array}\right) = 1,\]

以及 $\operatorname{rg}(A - 2E)^2 = \operatorname{rg}(0) = 0$, 从而对 $s \geq 2$ 也有 $\operatorname{rg}(A - 2E)^s = 0$.因此成立方程

\[
\begin{array}{l}
0 = \operatorname{rg}(A - 2E)^4 = n - n_1 = 0, \\
0 = \operatorname{rg}(A - 2E)^3 = k_{1, 4}, \\
0 = \operatorname{rg}(A - 2E)^2 = 2 k_{1, 4} + k_{1, 3}, \\
1 = \operatorname{rg}(A - 2E)^1 = 3 k_{1, 4} + 2 k_{1, 3} + k_{1, 2},
\end{array}\]

由此得到 $k_{1, 4} = k_{1, 3} = 0$, $k_{1, 2} = 1$, 以及

\[
k_{1, 1} = 4 - 2 k_{1, 2} = 2\text{.}\]

因此 $J = \operatorname{Diag}(J(2, 1), J(2, 1), J(2, 2))$, 即

\[
J = \left(\begin{array}{cccc} 2 & 0 & 0 & 0 \\ 0 & 2 & 0 & 0 \\ 0 & 0 & 2 & 0 \\ 0 & 0 & 1 & 2 \end{array}\right)\]

是 $A$ 的若尔当标准形.

最后我们再处理一个有效得多的方法, 借助它可以同时确定矩阵的若尔当标准形与一般标准形.做法如下: 从矩阵 $A \in K^{n \times n}$ 出发, 把它解释为 $n$ 维 $K$-向量空间 $V = K^n$ 的自同态 $f$, 从而可以在 $f$ 下把 $V$ 看作 $K[T]$-模.然后对 $A$ 构造 $V$ 作为 $K[T]$-模的典范有限展示, 并确定相应的初等因子.这样我们达到定理~\mtref{thm:6.4.2} 中给出的 $V$ 到单生成 $K[T]$-子模的分解, 从它可以得到一切进一步的分解, 特别也包括 $A$ 的标准形.

\begin{lemma}\label{lem:6.5.17}
对矩阵 $A \in K^{n \times n}$, 把 $K^n$ 看作由 $x \longmapsto A \cdot x$ 给出的自同态下的 $K[T]$-模.设

\[
\varphi \colon K[T]^n \longrightarrow K^n, \quad (p_1, \ldots, p_n) \longmapsto \sum_{i = 1}^{n} p_i(A) \cdot e_i,\]

是把 $K[T]^n$ 的典范 $K[T]$-基映到 $K^n$ 的典范 $K$-基 $e_1, \ldots, e_n$ 的 $K[T]$-线性映射; $\varphi$ 是满射.再设

\[
\psi \colon K[T]^n \longrightarrow K[T]^n, \quad P \longmapsto (TE - A) \cdot P,\]

是由矩阵 $TE - A \in K[T]^{n \times n}$ 给出的 $K[T]$-线性映射; $E \in K^{n \times n}$ 为单位矩阵.则序列

\[
K[T]^n \xrightarrow{\psi} K[T]^n \xrightarrow{\varphi} K^n \longrightarrow 0\]

是 $K^n$ 作为 $K[T]$-模的一个有限展示.
\end{lemma}

\begin{proof}
若 $e_i$ 在 $K[T]^n$ 中如 $K^n$ 中那样表示第 $i$ 个单位向量, 则对 $i = 1, \ldots, n$ 有

\[
\psi(e_i) = (TE - A)e_i = Te_i - Ae_i, \qquad \varphi(Te_i) = Ae_i, \qquad \varphi(Ae_i) = Ae_i,\]

从而

\[
\varphi \circ \psi(e_i) = \varphi(Te_i - Ae_i) = Ae_i - Ae_i = 0\text{.}\]

于是 $\varphi \circ \psi = 0$, 即 $\operatorname{im} \psi \subset \ker \varphi$.按同态定理, $\varphi$ 因此分解为满射的复合

\[
\varphi \colon K[T]^n \xrightarrow{\pi} K[T]^n/\operatorname{im} \psi \xrightarrow{\bar{\varphi}} K^n,\]

并且我们断言 $\bar{\varphi}$ 甚至是双射.由于 $K[T]^n$ 作为 $K$-向量空间由元素 $T^\nu e_i$($i = 1, \ldots, n$, $\nu \in \mathbb{N}$) 生成, 相应地, 像 $\pi(T^\nu e_i)$ 作为 $K$-向量空间生成商 $K[T]^n/\operatorname{im} \psi$.但对一切 $\nu \in \mathbb{N}$ 有

\[
T^{\nu + 1} e_j - T^\nu A e_j = \psi(T^\nu e_j) \in \operatorname{im} \psi, \qquad j = 1, \ldots, n,\]

因而

\[
T^{\nu + 1} e_j \in T^\nu \sum_{i = 1}^{n} K e_i + \operatorname{im} \psi, \qquad j = 1, \ldots, n\text{.}\]

由此对 $\nu \in \mathbb{N}$ 归纳得到

\[
T^\nu e_j \in \sum_{i = 1}^{n} K e_i + \operatorname{im} \psi, \qquad j = 1, \ldots, n,\]

从而

\[
K[T]^n = \sum_{i = 1}^{n} K e_i + \operatorname{im} \psi\text{.}\]

这意味着 $K[T]^n/\operatorname{im} \psi$ 作为 $K$-向量空间由 $n$ 个元素生成, 从而在 $K$ 上的维数 $\leq n$.按维数公式定理~\mtref{thm:2.1.10}, $\bar{\varphi}$ 于是必然是单射, 从而是双射, 如所断言.则 $\operatorname{im} \psi = \ker \pi = \ker \varphi$, 于是推出所给出的序列是 $K^n$ 作为 $K[T]$-模的有限展示.
\end{proof}

\begin{bigtheorem}\label{bigthm:6.5.18}
设 $f \colon V \longrightarrow V$ 是有限维 $K$-向量空间 $V$ 的自同态, 设 $A = A_{f, X, X} \in K^{n \times n}$ 是关于给定基 $X$ 描述 $f$ 的矩阵, 并设 $\alpha_1, \ldots, \alpha_s \in K[T]$ 满足 $\alpha_j \mid \alpha_{j + 1}$, 是矩阵 $TE - A \in K[T]^{n \times n}$ 的不可逆初等因子; 我们把 $\alpha_j$ 取作首一多项式.再考虑 $\alpha_j$ 的素因子分解 $\alpha_j = p_1^{n(1, j)} \ldots p_r^{n(r, j)}$, $j = 1, \ldots, s$, 其中 $p_1, \ldots, p_r \in K[T]$ 是两两不同的首一素多项式, $n(i, j) \geq 0$ 是指数.则:
\begin{enumerate}
\item[(i)] 在 $f$ 下把 $V$ 看作 $K[T]$-模, 则

\[
V \simeq \bigoplus_{j = 1}^{s} K[T]/\alpha_j K[T] \simeq \bigoplus_{i = 1}^{r} \bigoplus_{j = 1}^{s} K[T]/p_i^{n(i, j)} K[T],\]

并且后一个分解按定理~\mtref{bigthm:6.5.6} 给出 $V$ 到 $f$-循环线性子空间的分解.
\item[(ii)] $A$ 的一般标准形的方格是满足 $n(i, j) > 0$, $i = 1, \ldots, r$, $j = 1, \ldots, s$, 的素幂 $p_i^{n(i, j)}$ 的友矩阵.若一切 $p_i$ 都是线性的, 则由此利用引理~\mtref{lem:6.5.14} 得到若尔当标准形.
\item[(iii)] $f$ 的特征多项式 $\chi_f$ 与最小多项式 $p_f$ 算得为

\[
\chi_f = \alpha_1 \ldots \alpha_s, \qquad p_f = \alpha_s\text{.}\]
\end{enumerate}
\end{bigtheorem}

\begin{proof}
可以假设 $V = K^n$, 而 $f$ 是由 $x \longmapsto Ax$ 给出的 $V$ 的自同态.利用引理~\mtref{lem:6.5.17} 提供的 $V$ 的有限展示, 由定理~\mtref{thm:6.4.2} 得到 (i) 中的第一个分解, 并由中国剩余定理~\mtref{thm:6.4.1} 得到第二个, 其中定理~\mtref{thm:6.4.2} 中的自由部分消失, 因为 $V$ 是 $K[T]$-扭模.其余断言由定理~\mtref{thm:6.5.4} 以及引理~\mtref{lem:6.5.8} 和引理~\mtref{lem:6.5.10} 得出, 如定理~\mtref{bigthm:6.5.11} 的证明中那样.
\end{proof}

作为例子, 我们再次考虑矩阵

\[
A = \left(\begin{array}{cccc} 1 & 1 & 0 & 1 \\ 0 & 2 & 0 & 0 \\ -1 & 1 & 2 & 1 \\ -1 & 1 & 0 & 3 \end{array}\right) \in \mathbb{R}^{4 \times 4}\]

并先确定矩阵 $TE - A \in \mathbb{R}[T]^{4 \times 4}$ 的初等因子, 办法是应用定理~\mtref{lem:6.3.5} 证明中的方法:

\[
\left(\begin{array}{cccc} T - 1 & -1 & 0 & -1 \\ 0 & T - 2 & 0 & 0 \\ 1 & -1 & T - 2 & -1 \\ 1 & -1 & 0 & T - 3 \end{array}\right)
\longmapsto \left(\begin{array}{cccc} 1 & -1 & 0 & T - 3 \\ 0 & T - 2 & 0 & 0 \\ 1 & -1 & T - 2 & -1 \\ T - 1 & -1 & 0 & -1 \end{array}\right)\]

\[
\longmapsto \left(\begin{array}{cccc} 1 & -1 & 0 & T - 3 \\ 0 & T - 2 & 0 & 0 \\ 0 & 0 & T - 2 & -(T - 2) \\ 0 & T - 2 & 0 & -(T - 2)^2 \end{array}\right)
\longmapsto \left(\begin{array}{cccc} 1 & 0 & 0 & 0 \\ 0 & T - 2 & 0 & 0 \\ 0 & 0 & T - 2 & 0 \\ 0 & 0 & 0 & (T - 2)^2 \end{array}\right)\]

于是得到 $T - 2$, $T - 2$, $(T - 2)^2$ 为 $TE - A$ 的不可逆初等因子, 并按定理~\mtref{bigthm:6.5.18} 得到 $\chi_A = (T - 2)^4$ 为特征多项式, 以及 $p_A = (T - 2)^2$ 为 $A$ 的最小多项式.由于初等因子已经是素幂, 无需应用中国剩余定理, $A$ 的一般标准形与若尔当标准形得到为

\[
\left(\begin{array}{cccc} 2 & 0 & 0 & 0 \\ 0 & 2 & 0 & 0 \\ 0 & 0 & 0 & -4 \\ 0 & 0 & 1 & 4 \end{array}\right), \qquad \left(\begin{array}{cccc} 2 & 0 & 0 & 0 \\ 0 & 2 & 0 & 0 \\ 0 & 0 & 2 & 0 \\ 0 & 0 & 1 & 2 \end{array}\right)\text{.}\]

\aufgaben

下面设 $K$ 是域.

\noindent\textbf{习题 1.} 设 $A \in \mathbb{R}^{7 \times 7}$ 是特征多项式为

\[
\chi_A = (T^2 + 1)^2 (T - 2)(T^2 - 1) \in \mathbb{R}[T]\]

的矩阵.研究 $A$ 的一般标准形可能具有什么形状.

\medskip
\noindent\textbf{习题 2.} 证明: 两个矩阵 $A, B \in \mathbb{R}^{3 \times 3}$ 相似, 当且仅当它们的最小多项式以及特征多项式一致.(AT 357)

\medskip
\noindent\textbf{习题 3.} 对矩阵

\[
\left(\begin{array}{cccc} 2 & 2 & 0 & -3 \\ 1 & 1 & 0 & -1 \\ 1 & 2 & -1 & -1 \\ 1 & 2 & 0 & -2 \end{array}\right), \qquad \left(\begin{array}{cccc} 0 & 1 & 1 & -2 \\ 0 & 1 & 0 & 0 \\ 1 & 1 & 0 & -2 \\ -1 & 1 & 1 & -1 \end{array}\right), \qquad \left(\begin{array}{cccc} 3 & 0 & 1 & -2 \\ 4 & 1 & 1 & -3 \\ 0 & 0 & -1 & 0 \\ 4 & 0 & 2 & -3 \end{array}\right) \in \mathbb{R}^{4 \times 4}\]

分别计算特征多项式与最小多项式, 以及若存在则其若尔当标准形.

\medskip
\noindent\textbf{习题 4.} 设 $A, B \in K^{n \times n}$ 是最小多项式 $p_A, p_B \in K[T]$ 的矩阵, 带素因子分解 $p_A = p_1^{n_1} \cdot \ldots \cdot p_r^{n_r}$, 其中 $p_1, \ldots, p_r \in K[T]$ 是两两不同的首一素多项式.证明: $A$ 与 $B$ 相似, 当且仅当 $p_A = p_B$, 并且对 $i = 1, \ldots, r$ 和 $1 \leq k_i \leq n_i$ 有 $\operatorname{rg} p_i^{k_i}(A) = \operatorname{rg} p_i^{k_i}(B)$.

\medskip
\noindent\textbf{习题 5.} 考虑两两不同的首一素多项式 $p_1, \ldots, p_n \in K[T]$, 以及自然数 $1 \leq r_i \leq s_i$, $i = 1, \ldots, n$.是否存在有限维 $K$-向量空间 $V$ 的自同态 $f \colon V \longrightarrow V$, 使最小多项式 $p_f = p_1^{r_1} \cdot \ldots \cdot p_n^{r_n}$, 特征多项式 $\chi_f = p_1^{s_1} \cdot \ldots \cdot p_n^{s_n}$?

\medskip
\noindent\textbf{习题 6.} 设 $V$ 是有限维 $K$-向量空间, 带自同态 $f \colon V \longrightarrow V$, 并设 $U \subset V$ 是 $f$-不变的子空间.证明(AT 358):
\begin{enumerate}
\item[(i)] $f$ 诱导出自同态 $\bar{f} \colon V/U \longrightarrow V/U$;
\item[(ii)] 对 $\bar{f}$ 与 $f$ 的最小多项式有 $p_{\bar{f}} \mid p_f$;
\item[(iii)] 对 $f$, $f|_U$ 与 $\bar{f}$ 的特征多项式有 $\chi_f = \chi_{f|_U} \cdot \chi_{\bar{f}}$.
\end{enumerate}

\medskip
\noindent\textbf{习题 7.} 设 $V$ 是有限维 $K$-向量空间, 带自同态 $f \colon V \longrightarrow V$.设最小多项式 $p_f$ 是素多项式 $p \in K[T]$ 的幂, 比如 $p_f = p^r$ 且 $r > 0$.证明:
\begin{enumerate}
\item[(i)] 存在向量 $u \in V$ 使 $p^{r - 1}(f)(u) \neq 0$;
\item[(ii)] 若 $u \in V$ 如 (i) 所述, 且 $U \subset V$ 是由 $u$ 生成的 $f$-循环线性子空间, 则存在 $f$-不变的线性子空间 $U' \subset V$ 使 $V = U \oplus U'$.
\end{enumerate}

\medskip
\noindent\textbf{习题 8.} 设 $V$ 是在自同态 $f \colon V \longrightarrow V$ 下的有限维 $f$-循环 $K$-向量空间.证明 $V$ 的每个 $f$-不变的子空间 $U \subset V$ 也是 $f$-循环的.

\medskip
\noindent\textbf{习题 9.} 设 $V$ 是在自同态 $f \colon V \longrightarrow V$ 下的有限维 $f$-循环 $K$-向量空间, $u \in V$ 是生成元.证明:
\begin{enumerate}
\item[(i)] $V$ 上存在唯一确定的环结构, 其加法与 $V$ 上的向量空间加法一致, 其乘法满足 $(a f^i(u)) \cdot (b f^j(u)) = ab f^{i + j}(u)$, 其中 $a, b \in K$, $i, j \in \mathbb{N}$;
\item[(ii)] 在这个环结构下, $V$ 是域, 当且仅当 $f$ 的最小多项式是素的.
\end{enumerate}

